\documentclass[11pt]{article}
\usepackage[T1]{fontenc}
\usepackage{lmodern,amsmath,amssymb,amsthm}
\usepackage[margin=27mm]{geometry}
\usepackage{microtype,longtable,booktabs,array}
\usepackage[hyphens]{url}
\usepackage[hidelinks]{hyperref}
% The author's macros.
%
% some calligraphy
%
\newcommand\cA{{\mathcal A}}
\newcommand\cB{{\mathcal B}}
\newcommand\cC{{\mathcal C}}
\newcommand\cD{{\mathcal D}}
\newcommand\cE{{\mathcal E}}
\newcommand\cF{{\mathcal F}}
\newcommand\cG{{\mathcal G}}
\newcommand\cH{{\mathcal H}}
\newcommand\cI{{\mathcal I}}
\newcommand\cJ{{\mathcal J}}
\newcommand\cK{{\mathcal K}}
\newcommand\cL{{\mathcal L}}
\newcommand\cM{{\mathcal M}}
\newcommand\cN{{\mathcal N}}
\newcommand\cP{{\mathcal P}}
\newcommand\cQ{{\mathcal Q}}
\newcommand\cR{{\mathcal R}}
\newcommand\cS{{\mathcal S}}
\newcommand\cT{{\mathcal T}}
\newcommand\cU{{\mathcal U}}
\newcommand\cV{{\mathcal V}}

%
% fields and rings (and a semigroup)
%
\newcommand\cx{{\mathbb C}}% complexes
\newcommand\fld{{\mathbb F}}
\newcommand\flde{{\mathbb E}}
\newcommand\ints{{\mathbb Z}}
\newcommand\nn{{\mathbb N}}%non-negative integers
\newcommand\re{{\mathbb R}}%reals
\newcommand\rats{{\mathbb Q}}

\newcommand\Zu{{\mathbf u}}
\newcommand\Zv{{\mathbf v}}
\newcommand\Zw{{\mathbf w}}
\newcommand\Zy{{\mathbf y}}
\newcommand\Zz{{\mathbf z}}
\newcommand\Zx{{\mathbf x}}
\newcommand\Zj{{\mathbf j}}


%
% the really useful stuff
%
%\newcommand\opk[1]{\mathop{\hbox{\rm #1}}\nolimits} %% old version
\newcommand\opk[1]{\mathop{\mathrm{#1}}\nolimits}
\newcommand\comp[1]{{\mkern2mu\overline{\mkern-2mu#1}}}
\newcommand\diff{\mathbin{\mkern-1.5mu\setminus\mkern-1.5mu}}% for \setminus
\newcommand\gbin[2]{\genfrac{[}{]}{0pt}{}{#1}{#2}}
\newcommand\inv{^{-1}}
\newcommand\res{\mathbin{\mkern-2.0mu\restriction\mkern-2.0mu}}% ch2
\newcommand\sbs{\subseteq}
\newcommand\sps{\supseteq}
\newcommand\seq[3]{#1_{#2},\ldots,#1_{#3}}
\newcommand\supp{\opk{supp}}
\newcommand\im{\opk{im}}
\newcommand\row{\opk{row}}
\newcommand\pmat[1]{\begin{pmatrix} #1 \end{pmatrix}}
\newcommand\cprod{\mathbin{\square}}
%
% matrix theory
%
\newcommand\ip[2]{\langle#1,#2\rangle}
\newcommand\one{{\bf1}}
\newcommand\rk{\opk{rk}}
\newcommand\tr{\opk{tr}}
\newcommand\col{\opk{col}}
\newcommand\mat[3]{\mathrm{Mat}_{#1\times #2}(#3)}

%
% some group theory
%
\newcommand\aut[1]{{\rm Aut}(#1)}
\newcommand\fx[1]{{\rm fix}(#1)}% ch2
\newcommand\grp[1]{\langle #1\rangle}
\newcommand\nrml{\vartriangleleft}
\newcommand\nrmleq{\trianglelefteq}
\newcommand\alt[1]{\opk{Alt}(#1)}
\DeclareMathOperator{\Sym}{Sym}
% End of the author's macros.
\hypersetup{pdftitle={A spectral toolkit for CGSS},pdfauthor={}}
\newtheorem{theorem}{Theorem}[section]
\newtheorem{lemma}[theorem]{Lemma}
\newtheorem{corollary}[theorem]{Corollary}
\newtheorem{proposition}[theorem]{Proposition}
\theoremstyle{definition}
\newtheorem{definition}[theorem]{Definition}
\newtheorem{remark}[theorem]{Remark}
\DeclareMathOperator{\Spec}{Spec}
\DeclareMathOperator{\mult}{mult}
\DeclareMathOperator{\diag}{diag}
\DeclareMathOperator{\spanop}{span}
\makeatletter
\g@addto@macro\UrlBreaks{\do\a\do\b\do\c\do\d\do\e\do\f\do\g\do\h\do\i\do\j\do\k\do\l\do\m\do\n\do\o\do\p\do\q\do\r\do\s\do\t\do\u\do\v\do\w\do\x\do\y\do\z\do\A\do\B\do\C\do\D\do\E\do\F\do\G\do\H\do\I\do\J\do\K\do\L\do\M\do\N\do\O\do\P\do\Q\do\R\do\S\do\T\do\U\do\V\do\W\do\X\do\Y\do\Z}
\makeatother
\newcommand{\code}[1]{\ifmmode\text{\small\nolinkurl{#1}}\else{\small\nolinkurl{#1}}\fi}
\newcommand{\lean}[1]{\par\smallskip{\raggedright\small\noindent\textit{Proposed Lean: }\code{#1}.\par}}
\newcommand{\status}[1]{\par\smallskip{\raggedright\noindent\textit{Mathlib / new. }#1\par}}
\newcommand{\used}[1]{\par\smallskip\noindent\textit{Use. }#1\par}
\newcommand{\ms}[1]{\left[\!\left[#1\right]\!\right]}
\newcommand{\norm}[1]{\left\lVert#1\right\rVert}
\setlength{\emergencystretch}{2em}
\title{A spectral toolkit for CGSS\\\large Definitions, source catalogue, and ordinary proofs}
\author{}
\date{27 September 2026}
\begin{document}
\maketitle
\begin{abstract}
We give the spectral tools used in Sections 3--6 of Cameron, Goethals,
Seidel and Shult \cite{CGSS}, with multiplicities and all dimension and
nonemptiness hypotheses explicit. We will use the finite-dimensional real
spectral theorem from Mathlib as a foundation. The proofs below supply the
matrix and graph interfaces, symmetric Perron--Frobenius, the regular and
strongly regular graph calculations, and soundness theorems for exact
spectral certificates. The proposed names are the interface for the
structured proofs and the Lean code; not every one of them is a declaration
of the Lean formalization.
\end{abstract}

\section{Catalogue of existing results}\label{sec:catalogue}

We inspected Mathlib at commit \code{082e2d37} (4 September 2026), the
version pinned in \code{lake-manifest.json}. Paths in the following
catalogue are relative to \code{Mathlib/}. We record the actual
statements, rather than inferring their strength from declaration names
or introductory comments. ``New'' below means that the stated interface
needs a proof in this project, not that the mathematics is new.

\subsection{Mathlib foundations}

All the matrix statements in items M1--M5 specialize to finite real
matrices. A Lean matrix index has \code{Fintype} and, when matrix
multiplication or diagonal matrices require it, \code{DecidableEq}.
We will make these choices classically; no computation over real numbers
is part of a certificate checker.

\begin{description}\raggedright
\item[M1: the spectral theorem.]\label{cat:spectral}
\code{Analysis/Matrix/Spectrum.lean}, namespace
\code{Matrix.IsHermitian}, for a matrix over an \code{RCLike} field and
\code{hA : A.IsHermitian}, provides
\code{eigenvalues}\textsubscript{\code{0}} (a decreasing real list indexed by
\code{Fin (Fintype.card U)}), \code{eigenvalues} (the same list reindexed
by \code{U}), \code{eigenvalues}\textsubscript{\code{0}}\code{_antitone},
\code{eigenvectorBasis}, \code{mulVec_eigenvectorBasis}, and
\code{spectral_theorem}. The last asserts unitary diagonalization.
The same namespace contains \code{charpoly_eq},
\code{roots_charpoly_eq_eigenvalues},
\code{roots_charpoly_eq_eigenvalues}\textsubscript{\code{0}}, \code{splits_charpoly},
\code{det_eq_prod_eigenvalues}, \code{trace_eq_sum_eigenvalues}, and
\code{rank_eq_card_non_zero_eigs}. These give, respectively, the product
of linear factors, the multiset of roots, splitting, product, sum, and
number of nonzero eigenvalues. Also,
\code{eigenvalues_eq_eigenvalues_iff} equates equality of the reindexed
eigenvalue functions with equality of characteristic polynomials for two
Hermitian matrices on the \emph{same} index type;
\code{spectrum_real_eq_range_eigenvalues} identifies the ordinary
set-valued algebra spectrum. In particular the root theorem's full name
is \code{Matrix.IsHermitian.roots_charpoly_eq_eigenvalues}.

\item[M2: eigenspaces and coordinate changes.]\label{cat:eigenspace}
\code{Analysis/InnerProductSpace/Spectrum.lean} contains
\code{LinearMap.IsSymmetric.card_filter_eigenvalues_eq}: if $T$ is
symmetric on a finite-dimensional inner product space over an
\code{RCLike} field and $h_n:\dim V=n$, the number of indices $i$ with
$\operatorname{eigenvalues}(T,h_n)(i)=\mu$ is
$\dim\ker(T-\mu I)$ (with the real eigenvalue cast into the field).
The name is \emph{card\_filter\_eigenvalues\_eq}.
\code{Analysis/Matrix/Hermitian.lean} contains
\code{Matrix.isSymmetric_toEuclideanLin_iff}:
$A.\mathrm{toEuclideanLin}$ is symmetric iff $A$ is Hermitian.
\code{Matrix.isHermitian_iff_isSymmetric} is its deprecated reverse
alias. \code{Matrix.toEuclideanLin} is in
\code{Analysis/InnerProductSpace/PiL2.lean}.
\code{Module.End.eigenspace} and \code{Module.End.mem_eigenspace_iff}
are in \code{LinearAlgebra/Eigenspace/Basic.lean}.
\code{LinearAlgebra/Eigenspace/Zero.lean} supplies
\code{LinearMap.finrank_eigenspace_le} and
\code{LinearMap.finrank_maxGenEigenspace_eq}: over a field and a finite
module, geometric multiplicity is at most characteristic-polynomial
root multiplicity, and maximal generalized-eigenspace dimension equals
that multiplicity. The equality for a symmetric matrix's ordinary
eigenspace in \code{toLin'} coordinates is a new bridge, proved below.
\code{Polynomial.count_roots}, in
\code{Algebra/Polynomial/Roots.lean}, identifies root count with
\code{rootMultiplicity}.

\item[M3: positivity.]\label{cat:positive}
\code{LinearAlgebra/Matrix/PosDef.lean} contains
\code{Matrix.posSemidef_iff_dotProduct_mulVec}: positivity is Hermitian
symmetry together with nonnegativity of every quadratic form. It also
contains \code{Matrix.PosSemidef.submatrix}.
\code{Analysis/Matrix/PosDef.lean} contains
\code{Matrix.IsHermitian.posSemidef_iff_eigenvalues_nonneg},
\code{Matrix.PosSemidef.eigenvalues_nonneg},
\code{Matrix.IsHermitian.posDef_iff_eigenvalues_pos}, and
\code{Matrix.PosSemidef.det_nonneg}, for finite Hermitian matrices over
an \code{RCLike} field.
\code{Matrix.PosSemidef.dotProduct_mulVec_zero_iff}, in
\code{Analysis/Matrix/Order.lean}, says that $\Zx^*A\Zx=0$ iff
$A\Zx=0$ when $A$ is positive semidefinite.
\code{Matrix.posSemidef_iff_eq_sum_vecMulVec}, in
\code{Analysis/InnerProductSpace/Positive.lean}, expresses a positive
semidefinite matrix as a finite sum of $\Zv\Zv^*$; it does not itself
specify a representation of dimension equal to the rank.
\code{Matrix.rank_transpose_mul_self} and
\code{Matrix.rank_self_mul_transpose}, in
\code{LinearAlgebra/Matrix/Rank.lean}, give
$\rk(B^TB)=\rk B=\rk(BB^T)$ over a linearly ordered field.

\item[M4: characteristic polynomials and trace.]\label{cat:charpoly}
\code{LinearAlgebra/Matrix/Charpoly/Basic.lean}, over a commutative ring,
contains \code{Matrix.charpoly_mul_comm'} for rectangular matrices:
$t^q\chi_{MN}=t^p\chi_{NM}$ for $M$ of size $p\times q$;
\code{Matrix.charpoly_mul_comm} is the square version, and
\code{Matrix.charpoly_diagonal} is the product of linear factors for a
diagonal matrix. \code{Polynomial.roots_multiset_prod_X_sub_C}, in
\code{Algebra/Polynomial/Roots.lean}, reads the roots of such a product.
\code{Matrix.charpoly_units_conj} and
\code{Matrix.charpoly_units_conj'} give conjugation invariance by a
matrix unit. \code{Matrix.charpoly_map} commutes characteristic
polynomials with a ring homomorphism.
\code{LinearMap.IsProj.trace}, in \code{LinearAlgebra/Trace.lean},
equates the trace of a projection with the dimension of its image;
the general ring statement assumes finite free image and kernel.
\code{LinearMap.isProj_range_iff_isIdempotentElem}, in
\code{LinearAlgebra/Projection.lean}, supplies its idempotent hypothesis.
\code{Matrix.det_vandermonde_ne_zero_iff} and
\code{Matrix.eq_zero_of_forall_pow_sum_mul_pow_eq_zero}, in
\code{LinearAlgebra/Vandermonde.lean}, give invertibility and uniqueness
of moments at distinct nodes over a domain.
\code{isInteger_of_is_root_of_monic}, in
\code{RingTheory/Polynomial/RationalRoot.lean}, says that a root in a
fraction field of a monic polynomial over a UFD lies in the base ring.
We will use the specialization $\ints\sbs\rats$.

\item[M5: Rayleigh quotients.]\label{cat:rayleigh}
\code{Analysis/InnerProductSpace/Rayleigh.lean} contains
\code{LinearMap.IsSymmetric.hasEigenvalue_iSup_of_finiteDimensional}
and \code{LinearMap.IsSymmetric.hasEigenvalue_iInf_of_finiteDimensional}:
the supremum and infimum of the Rayleigh quotient on \emph{nonzero}
vectors are eigenvalues, assuming finite dimension and a nontrivial
space. It also contains
\code{IsSelfAdjoint.hasEigenvector_of_isMaxOn} for continuous linear
maps. We will prove the matrix maximum and its equality condition
directly from M1. Mathlib has an irreducibility definition,
\code{Matrix.IsIrreducible}, in
\code{LinearAlgebra/Matrix/Irreducible/Defs.lean}, but no
Perron--Frobenius theorem and no Cauchy interlacing theorem found by our
source search. Our support-graph formulation includes the one-vertex
zero matrix without requiring a positive-length loop.

\item[M6: graph matrices.]\label{cat:graph}
\code{Combinatorics/SimpleGraph/AdjMatrix.lean} contains
\code{SimpleGraph.adjMatrix_pow_apply_eq_card_walk}: over a semiring,
$(A^j)(u,v)$ is the cardinality of walks of length $j$ from $u$ to $v$,
cast to that semiring. It also contains
\code{SimpleGraph.trace_adjMatrix},
\code{SimpleGraph.adjMatrix_mul_self_apply_self} (the diagonal of $A^2$
is degree), \code{SimpleGraph.adjMatrix_mulVec_const_apply_of_regular},
and \code{SimpleGraph.Iso.reindex_adjMatrix}. For the complement,
\code{SimpleGraph.compl_adjMatrix_eq_adjMatrix_compl} identifies
$A(\comp X)$ with Mathlib's \code{Matrix.compl} of $A(X)$, and
\code{SimpleGraph.one_add_adjMatrix_add_compl_adjMatrix_eq_of_one} gives
$I+A(X)+A(X).\mathrm{compl}=J$. The structure \code{Matrix.IsAdjMatrix}
bundles the three guards zero-one entries, symmetry and zero diagonal;
\code{Matrix.IsAdjMatrix.toGraph} is the graph with adjacency $A(u,v)=1$,
and \code{Matrix.IsAdjMatrix.adjMatrix_toGraph_eq} says that its
adjacency matrix is $A$.
\code{Combinatorics/SimpleGraph/LapMatrix.lean} contains
\code{SimpleGraph.isHermitian_adjMatrix}, \code{SimpleGraph.lapMatrix},
\code{SimpleGraph.posSemidef_lapMatrix},
\code{SimpleGraph.lapMatrix_mulVec_eq_zero_iff_forall_reachable}, and
\code{SimpleGraph.card_connectedComponent_eq_finrank_ker_toLin'_lapMatrix}.
The last two are real statements for finite graphs with decidable
adjacency, and identify the Laplacian kernel with functions constant on
components. \code{SimpleGraph.incMatrix_mul_transpose}, in
\code{Combinatorics/SimpleGraph/IncMatrix.lean}, gives the entries of
$MM^T=D+A$ over a semiring. Its columns run over \emph{all} of \code{Sym2 U}; columns
which are not edges are zero. We will restrict to actual edges before
forming $M^TM$ for a line graph.
\code{SimpleGraph.lineGraph}, in
\code{Combinatorics/SimpleGraph/LineGraph.lean}, has vertex type
\code{X.edgeSet}; two distinct edges are adjacent when they share
a vertex.

\item[M7: strongly regular graphs.]\label{cat:srg}
\code{Combinatorics/SimpleGraph/StronglyRegular.lean} contains
\code{SimpleGraph.IsSRGWith.matrix_eq} over a semiring,
\code{SimpleGraph.IsSRGWith.param_eq} for $n>0$, and
\code{SimpleGraph.IsSRGWith.compl}. The matrix identity uses the
complement adjacency matrix; the parameter formulas use natural-number
subtraction. Complete and empty graphs are admitted by
\code{IsSRGWith}. The three-eigenvalue theorem below adds connectedness
and $0<k<n-1$, and performs the arithmetic in $\re$ or $\ints$.

\item[M8: auxiliary matrix and counting facts.]\label{cat:auxiliary}
\code{Matrix.det_one_add_replicateCol_mul_replicateRow}, in
\code{LinearAlgebra/Matrix/SchurComplement.lean}, proves
$\det(I+\Zu\Zv^T)=1+\Zv^T\Zu$ over a commutative ring.
\code{Sym2.card}, in \code{Data/Sym/Card.lean}, gives
$|\operatorname{Sym2}(U)|=\binom{n+1}{2}$ for a finite type of
cardinality $n$. This supplies the count in the symmetric-matrix
dimension argument.
\code{Matrix.IsHadamard}, in
\code{LinearAlgebra/Matrix/HadamardMatrix.lean}, is defined over a
star semiring by unitary entries and both equations $HH^*=nI$ and
$H^*H=nI$. Over $\re$ this means entries $\pm1$ and the two
orthogonality equations. Our Hadamard theorem explicitly derives
the second equation from the first when $n>0$, so its hypotheses
produce this existing predicate.
\end{description}

\subsection{The new interfaces and their uses}

We will use \code{Matrix.spectrumMultiset}, rather than redefine
Mathlib's set-valued \code{spectrum}, and \code{SimpleGraph.Cospectral}
for cospectral graphs. The remaining
new interfaces are multiplicity and eigenspace dimension, polynomial
spectral mapping with multiplicities, ordered extremes, their positivity
bridges, symmetric Perron--Frobenius, the signless Laplacian kernel,
line-graph spectral transfer, strongly regular spectral idempotents, and
the remaining certificate wrappers. Each result below has its Lean name, status,
and place of use. The names for Lemmas 5.1--5.6 and Smith's theorem follow
\code{cgss/notation.md}, Section 10.

\section{Definitions and the spectral theorem}\label{sec:definitions}

We will use functional notation for entries: $A(u,v)$ and $\Zv(u)$.
Vectors are columns, $I$ is the identity matrix, and $J$ is the all-ones
matrix; $\Zj$ is the all-ones vector. For a finite set $U$ of cardinality
$n$, $\re^U$ has inner product
$\ip{\Zx}{\Zy}=\sum_{u\in U}\Zx(u)\Zy(u)$.
The \emph{rank} of a matrix is the dimension of the image of its
associated linear map. All graphs are finite and simple. Matrices in
spectral statements are real and symmetric unless another hypothesis is
specified. An empty matrix has characteristic polynomial $1$, determinant
$1$, rank $0$, and empty spectrum.

\begin{definition}\label{def:spectrum}
For a real square matrix $A$, put
\[
 \chi_A(t)=\det(tI-A),\qquad
 \Spec(A)=\chi_A.\mathrm{roots},\qquad
 \mult_A(\mu)=\operatorname{count}_{\Spec(A)}(\mu).
\]
Here roots form a multiset over $\re$, with algebraic multiplicity.
We will call $\ker(A-\mu I)$ the \emph{eigenspace} $V_A(\mu)$; an
eigenvector is a nonzero member. The definition of $\Spec$ makes sense
without symmetry, but its cardinality can then be smaller than $n$.
\lean{Matrix.spectrumMultiset, Matrix.eigenvalueMultiplicity}
\status{Implemented in \code{LineGraphs/Spectral/Basic.lean}, using \code{Matrix.charpoly},
\code{Polynomial.roots}, and \code{Module.End.eigenspace} (M1, M2).}
\end{definition}

\begin{theorem}[Coordinates, roots, and multiplicity]\label{thm:multiplicity}
Let $A$ be a real symmetric matrix on a finite set $U$ of cardinality $n$.
There is an orthonormal basis $(\Zv_i)_{i\in U}$ with
$A\Zv_i=\lambda_i\Zv_i$, where
$\lambda_i=\code{Matrix.IsHermitian.eigenvalues}(A)(i)$, and
\[
 \chi_A(t)=\prod_{i\in U}(t-\lambda_i),\quad
 \Spec(A)=\ms{\lambda_i:i\in U},\quad
 \mult_A(\mu)=\dim V_A(\mu)\quad(\mu\in\re).
\]
Thus $|\Spec(A)|=n$; membership in $\Spec(A)$, existence of a nonzero
$\mu$-eigenvector, and membership in Mathlib's set-valued
$\operatorname{spectrum}_{\re}(A)$ are equivalent. Multiplicity also
equals \code{A.charpoly.rootMultiplicity mu}.
More generally, if $(\Zu_j)_{j\in I}$ is any orthonormal basis of $\re^U$
with $A\Zu_j=\mu_j\Zu_j$, then $\Spec(A)=\ms{\mu_j:j\in I}$.
There is also a sorted orthonormal eigenbasis $(\Zw_i)_{i\in\operatorname{Fin}(n)}$
with $A\Zw_i=\nu_i\Zw_i$, where
$\nu_i=\code{hA.eigenvalues}\textsubscript{0}(i)$ and $\nu_i\ge\nu_j$ for $i\le j$.
Its eigenvalues list the spectrum in decreasing order:
$\Spec(A)=\ms{\nu_i:i\in\operatorname{Fin}(n)}$.
\end{theorem}
\begin{proof}
Over $\re$, conjugation is trivial, so symmetry is exactly the Hermitian
hypothesis of M1. Apply its orthonormal eigenbasis theorem, and form the
matrix $Q$ whose columns are these vectors. Orthonormality gives
$Q^TQ=QQ^T=I$, and $AQ=Q\diag(\lambda_i)$. Thus
$A=Q\diag(\lambda_i)Q^T$. Taking the determinant of
$tI-A=Q(tI-\diag(\lambda_i))Q^T$ gives the displayed factorization;
$\det Q\det Q^T=1$. Each factor has one root with multiplicity one,
and multiplying factors adds their root multisets. This proves the
multiset assertion, including repeated eigenvalues. For an arbitrary
orthonormal eigenbasis $(\Zu_j)_{j\in I}$, let $P$ have these columns and
$D=\diag(\mu_j)$. Both $A$ and $PDP^T=\sum_j\mu_j\Zu_j\Zu_j^T$ act on each
$\Zu_k$ by $\mu_k$, so $A=PDP^T$. Since $P^TP=I$ and $|I|=n$,
\code{Matrix.charpoly_mul_comm'} gives
$t^n\chi_A=t^n\chi_{DP^TP}=t^n\chi_D$, and
\code{Matrix.charpoly_diagonal} and
\code{Polynomial.roots_multiset_prod_X_sub_C} give the multiset.

Write $\Zx=\sum_i b_i\Zv_i$ in this basis. The equation
$A\Zx=\mu\Zx$ is equivalent to
$(\lambda_i-\mu)b_i=0$ for every $i$, by uniqueness of basis
coefficients. Thus the vectors $\Zv_i$ with $\lambda_i=\mu$ form a
basis of $V_A(\mu)$: they are independent as a subset of a basis, and
the preceding coefficient equation proves spanning. Their number is
the multiset count. A positive count is equivalent to a nonzero
eigenspace. Also, $\mu I-A$ fails to be invertible exactly when its
determinant is zero, equivalently when one factor $\mu-\lambda_i$
vanishes; this is the definition of the algebra spectrum. Finally
\code{Polynomial.count_roots} identifies the root count with
\code{rootMultiplicity}.

For the Lean coordinate bridge, send $\Zx:U\to\re$ to
$\code{WithLp.toLp}\ 2\ \Zx$. This is a linear equivalence with
$\code{EuclideanSpace}\ \re\ U$, with inverse taking the underlying
function. The formula for \code{toEuclideanLin} says that this map
intertwines it with \code{A.toLin'}. It thus restricts to inverse
linear maps on the two eigenspaces, preserving their finite dimensions.
Apply \code{LinearMap.IsSymmetric.card_filter_eigenvalues_eq} to
$A.\mathrm{toEuclideanLin}$ and $\dim\re^U=n$; reindexing the finite
count gives the asserted multiplicity in either coordinate convention.
For the sorted basis, let
$e:\operatorname{Fin}(n)\simeq U$ be
\code{Fintype.equivOfCardEq (Fintype.card_fin n)} and set
$\Zw_i=\code{hA.eigenvectorBasis}(e(i))$.
This is \code{hA.eigenvectorBasis.reindex e.symm}.
The definition of \code{hA.eigenvalues} gives
$\lambda_{e(i)}=\nu_i$; apply \code{mulVec_eigenvectorBasis} at $e(i)$.
Finally \code{eigenvalues}\textsubscript{0}\code{_antitone} gives the ordering, and
\code{roots_charpoly_eq_eigenvalues}\textsubscript{0} gives the sorted multiset.
If $U$ is empty, all bases and multisets above
are empty and both dimensions are zero.
\end{proof}
\lean{Matrix.IsHermitian.spectrumMultiset_eq_eigenvalues, Matrix.IsHermitian.eigenvalueMultiplicity_eq_finrank, Matrix.IsHermitian.mem_spectrumMultiset_iff, Matrix.IsHermitian.sortedEigenvectorBasis, Matrix.spectrumMultiset_eq_of_orthonormal_eigenbasis}
\status{Mathlib: M1, M2 and M4. The root-multiset wrapper is implemented in
\code{LineGraphs/Spectral/Basic.lean}; the \code{toLin'} bridge and sorted-basis
wrapper are proposed here. The arbitrary-basis assertion needs no symmetry
hypothesis; it is proved in Lean in \code{LineGraphs/Spectral/Multiplicity.lean}.}
\used{All spectral statements; specifically [CGSS] Lemmas 5.1--5.6,
SR3--SR4, and the rank arguments of Sections 4 and 6.}

\begin{definition}\label{def:extremes}
For a real symmetric matrix $A$ on $n>0$ indices, define
\[
 \lambda_{\max}(A)=\code{hA.eigenvalues}\textsubscript{0}(0),\qquad
 \lambda_{\min}(A)=\code{hA.eigenvalues}\textsubscript{0}(n-1).
\]
The indices carry their proofs of being less than $n$. These are exactly
\code{LineGraphs/Spectral/Basic.lean}'s definitions. They require
\code{Nonempty U} and \code{hA : A.IsHermitian}; there is no extreme
for the empty matrix.
\lean{Matrix.IsHermitian.largestEigenvalue, Matrix.IsHermitian.leastEigenvalue}
\status{Implemented first- and last-entry wrappers in \code{LineGraphs/Spectral/Basic.lean}.}
\end{definition}

\begin{lemma}[Extremal entries and multiset bounds]\label{lem:extremes}
Let $A$ be real symmetric on a nonempty finite index set of size $n$.
Then $\lambda_{\max}(A),\lambda_{\min}(A)\in\Spec(A)$, and
\[
 \lambda_{\min}(A)\le\mu\le\lambda_{\max}(A)
 \qquad(\mu\in\Spec(A)).
\]
Thus they are the maximum and minimum of the multiset and equal
the first and last entries of the decreasing list.
\end{lemma}
\begin{proof}
Apply \code{Matrix.IsHermitian.roots_charpoly_eq_eigenvalues}\textsubscript{0}.
Each endpoint index supplies membership, and each member has the form
$\nu_i$ for some $0\le i\le n-1$. Applying
\code{Matrix.IsHermitian.eigenvalues}\textsubscript{0}\code{_antitone}
to $0\le i$ and $i\le n-1$ gives both bounds. Membership and these
bounds are precisely the maximum and minimum properties.
\end{proof}
\lean{Matrix.IsHermitian.largestEigenvalue_mem_spectrumMultiset, Matrix.IsHermitian.leastEigenvalue_mem_spectrumMultiset, Matrix.IsHermitian.le_largestEigenvalue_of_mem, Matrix.IsHermitian.leastEigenvalue_le_of_mem}
\status{New ordering bridges from M1, for the existing Lean definitions.}
\used{Rayleigh bounds, interlacing and every extremal-eigenvalue comparison.}

\begin{definition}\label{def:graph}
The spectrum and multiplicity of a graph $X$ are those of its adjacency
matrix $A(X)$. Two graphs, possibly on different vertex types, are
\emph{cospectral} when their real adjacency characteristic polynomials
are equal. The least and largest graph eigenvalues require a nonempty
vertex set. The \emph{Seidel matrix} is
\[
 S(X)=J-I-2A(X).
\]
It is symmetric, has diagonal zero, and off-diagonal entry $-1$ on an
edge and $1$ on a nonedge.
\lean{SimpleGraph.spectrumMultiset, SimpleGraph.eigenvalueMultiplicity, SimpleGraph.Cospectral, SimpleGraph.largestEigenvalue, SimpleGraph.leastEigenvalue, SimpleGraph.seidelMatrix}
\status{Implemented in \code{LineGraphs/Spectral/Basic.lean} over
Mathlib's \code{SimpleGraph.adjMatrix}.}
\end{definition}

\begin{proposition}[Equivalent spectra]\label{prop:cospectral}
For real symmetric matrices $A,B$, even on different finite index sets,
$\chi_A=\chi_B$ iff $\Spec(A)=\Spec(B)$ iff
$\mult_A(\mu)=\mult_B(\mu)$ for all real $\mu$.
Thus cospectral graphs have the same number of vertices. Graph
isomorphisms preserve the spectrum.
\end{proposition}
\begin{proof}
Equality of polynomials gives equality of their root multisets.
Conversely, Theorem~\ref{thm:multiplicity} expresses each polynomial
as the product of $t-\mu$ over its multiset, with leading coefficient
one, so equal multisets give equal polynomials. Two finite multisets
are equal iff all their counts are equal: remove one occurrence of an
element and argue by induction on the size, with the empty multiset
as base case. Their common size gives the common matrix order.
For an isomorphism, relabel the coordinate basis by its vertex
bijection. The two adjacency linear maps then have the same matrix
in corresponding bases. In the determinant expansion, the relabelling
bijects the permutations, preserves their signs, and preserves every
entry product. Thus their characteristic polynomials are equal.
\end{proof}
\lean{Matrix.IsHermitian.spectrumMultiset_eq_iff_charpoly_eq, SimpleGraph.cospectral_iff_spectrumMultiset_eq, SimpleGraph.Iso.cospectral}
\status{The isomorphism wrapper is implemented in
\code{LineGraphs/Spectral/Basic.lean}. The equivalence wrappers are proposed
from M1 and \code{SimpleGraph.Iso.reindex_adjMatrix} (M6).}
\used{[CGSS] Section 5, C13 and C15.}

\section{Polynomial calculus, traces, and positive semidefiniteness}\label{sec:calculus}

\begin{theorem}[Polynomial calculus and spectral idempotents]\label{thm:calculus}
Let $A$ be real symmetric of order $n$. List its distinct eigenvalues as
$\theta_0>\cdots>\theta_d$. For $n=0$ the list is empty. There are
unique symmetric matrices $E_r$ that are the identity on $V_A(\theta_r)$
and zero on every other eigenspace. They satisfy
\[
 \sum_rE_r=I,\quad E_rE_s=\begin{cases}E_r&r=s,\\0&r\ne s,\end{cases}
 \qquad A=\sum_r\theta_rE_r.
\]
For every real polynomial $f$,
\[
 f(A)=\sum_r f(\theta_r)E_r,\quad
 \Spec(f(A))=\operatorname{map}(f,\Spec(A)),\quad
 \mult_{f(A)}(b)=\sum_{r:f(\theta_r)=b}\mult_A(\theta_r).
\]
In particular $f(A)=0$ iff $f(\theta_r)=0$ for every $r$, and
\[
 E_r=\prod_{s\ne r}\frac{A-\theta_sI}{\theta_r-\theta_s},\quad
 \tr E_r=\rk E_r=\mult_A(\theta_r).
\]
\end{theorem}
\begin{proof}
In the orthonormal eigenbasis from Theorem~\ref{thm:multiplicity},
define $E_r$ to have diagonal entry $1$ at those indices with
$\lambda_i=\theta_r$ and $0$ elsewhere, then conjugate by $Q$.
These matrices are symmetric. Their actions on each basis vector
prove the sum, product, and decomposition identities, as well as
uniqueness. For $j=0$, $A^0=I$ acts by $1=\lambda_i^0$ even when
$\lambda_i=0$. If $A^j\Zv_i=\lambda_i^j\Zv_i$, multiplication by $A$
proves the assertion for $j+1$. Taking linear combinations of these
power identities proves the formula for $f(A)$.
Thus $f(A)=Q\diag(f(\lambda_i))Q^T$. Taking its characteristic
polynomial as before proves the multiset formula; counting equal
images gives the multiplicity formula, with collisions added.
The zero-matrix equivalence follows by applying the matrix to each
nonzero basis vector. The polynomial in the displayed product has
value $1$ at $\theta_r$ and $0$ at every other $\theta_s$, since
distinctness makes all denominators nonzero. It thus acts as
$E_r$ on the eigenbasis. Its rank is the number of diagonal ones.
For its trace, expand
$\sum_u E_r(u,u)=\sum_{i:\lambda_i=\theta_r}\sum_u\Zv_i(u)^2$;
every inner sum is $1$ by normalization. This gives the same count.
\end{proof}
\lean{Matrix.IsHermitian.spectralIdempotent, Matrix.IsHermitian.aeval_eq_sum_spectralIdempotent, Matrix.IsHermitian.spectrumMultiset_aeval, Matrix.IsHermitian.aeval_eq_zero_iff, Matrix.IsHermitian.trace_spectralIdempotent}
\status{New construction, multiset calculus and idempotent algebra from M1.}
\used{SR3, [CGSS] Theorems 3.7, 4.4, 5.8, Section 6, and certificates.}

\begin{theorem}[Trace, determinant, rank, and moments]\label{thm:trace}
If $A$ is real symmetric with eigenvalues $\lambda_1,\ldots,\lambda_n$
listed with multiplicity, then, for every integer $j\ge0$,
\[
 \tr A^j=\sum_{i=1}^n\lambda_i^j,\qquad
 \det A=\prod_{i=1}^n\lambda_i,\qquad
 \rk A=n-\mult_A(0).
\]
For any real polynomial $f$,
$\tr f(A)=\sum_i f(\lambda_i)$ and
$\rk f(A)=\sum_{r:f(\theta_r)\ne0}\mult_A(\theta_r)$.
If $E$ is any real idempotent square matrix, even if not symmetric,
$\tr E=\rk E$ and this number is a nonnegative integer.
\end{theorem}
\begin{proof}
Use the diagonalization of $f(A)$ in Theorem~\ref{thm:calculus}.
For rectangular products for which both traces are defined,
$\tr(PQ)=\sum_{u,v}P(u,v)Q(v,u)=\tr(QP)$, by interchanging the
finite sums. Thus trace is unchanged by conjugation. Determinant
is unchanged by conjugation by multiplicativity; rank is unchanged
because an invertible map bijects both domain and image. Diagonal
trace, determinant, and rank are respectively the sum, product, and
number of nonzero diagonal entries. This proves all assertions for $A$.
For an idempotent $E$, every vector has the decomposition
$\Zx=E\Zx+(\Zx-E\Zx)$ into image and kernel. Their intersection is
zero: if $\Zy=E\Zz$ and $E\Zy=0$, then
$\Zy=E^2\Zz=E\Zy=0$. A basis of the image followed by a basis of
the kernel thus gives a matrix $\diag(I_{\rk E},0)$ for $E$.
Trace invariance proves the last assertion.
\end{proof}
\lean{Matrix.IsHermitian.trace_pow_eq_sum_spectrumMultiset, Matrix.IsHermitian.rank_aeval_eq_sum_multiplicity, Matrix.trace_eq_rank_of_idempotent}
\status{Mathlib: determinant, rank and first trace in M1; idempotent
trace in M4. The multiset power trace is implemented in
\code{LineGraphs/Spectral/Basic.lean}; the polynomial rank and matrix
idempotent wrappers are proposed here.}
\used{SP2, SR3--SR4, EQ1--EQ3, [CGSS] Theorem 3.7, Section 6, C13--C15.}

\begin{theorem}[Positivity and shifts]\label{thm:psd}
Let $A$ be real symmetric. It is positive semidefinite iff every member
of $\Spec(A)$ is nonnegative, and positive definite iff every member is
positive. For any $b\in\re$,
\[
 A-bI\succeq0\ \Longleftrightarrow\ (\forall\mu\in\Spec(A))\ b\le\mu,
 \qquad
 bI-A\succeq0\ \Longleftrightarrow\ (\forall\mu\in\Spec(A))\ \mu\le b.
\]
If $A\succeq0$, then $\Zx^TA\Zx=0$ iff $A\Zx=0$.
If $A\succeq0$, every principal submatrix is positive semidefinite and has
nonnegative determinant. For positive semidefinite real matrices $A,B$ of the same
order, $\tr(AB)\ge0$.
\end{theorem}
\begin{proof}
Write $\Zx=\sum_i b_i\Zv_i$. Then
$\Zx^TA\Zx=\sum_i\lambda_i b_i^2$ and
$\norm{\Zx}^2=\sum_i b_i^2$. If all $\lambda_i$ are nonnegative,
the first sum is nonnegative. Conversely, substitute each $\Zv_i$.
If all are positive, every nonzero vector has a nonzero coefficient,
so the sum is positive; conversely substitute each $\Zv_i$ again.
The shifted matrices act on $\Zv_i$ by $\lambda_i-b$ and $b-\lambda_i$,
which proves both equivalences. In a sum of nonnegative terms,
vanishing forces every term to vanish; thus
$\sum_i\lambda_i b_i^2=0$ iff every $\lambda_i b_i=0$, which is
$A\Zx=0$. To test a principal submatrix's form, extend a vector by
zeros outside its index subset. Its form equals that for $A$, so is
nonnegative. The determinant assertion follows from
Theorem~\ref{thm:trace}, since it is a product of nonnegative eigenvalues.
Finally diagonalize $A=\sum_i\lambda_i\Zv_i\Zv_i^T$. Finite summation
and the trace formula give
$\tr(AB)=\sum_i\lambda_i\Zv_i^TB\Zv_i\ge0$.
\end{proof}
\lean{Matrix.IsHermitian.posSemidef_iff_spectrumMultiset_nonneg, Matrix.IsHermitian.posSemidef_sub_smul_one_iff, Matrix.IsHermitian.posSemidef_smul_one_sub_iff}
\status{Mathlib: M3. New: multiset and shift wrappers, and the trace
inequality $\tr(AB)\ge0$.}
\used{LA4, HC3, [CGSS] Sections 3--4 and 6, the standing least-eigenvalue hypothesis.}

\begin{corollary}[The least-eigenvalue bridge]\label{cor:minus-two}
For every finite graph $X$ with adjacency matrix $A$,
\[
 (2I+A)\succeq0\ \Longleftrightarrow\
 (\forall\mu\in\Spec(X))\ -2\le\mu,
 \qquad -2\in\Spec(X)\ \Longleftrightarrow\ \det(2I+A)=0.
\]
If $X$ has a vertex, the first right-hand side is
$\lambda_{\min}(X)\ge-2$. Its least eigenvalue equals $-2$ iff both
the positive semidefiniteness and determinant-zero conditions hold.
Also $2I-A\succeq0$ iff $\lambda_{\max}(X)\le2$ for a nonempty graph.
\end{corollary}
\begin{proof}
Adjacency is symmetric. Theorem~\ref{thm:psd} with $b=-2$ gives the
first equivalence. Theorem~\ref{thm:calculus} gives the eigenvalues
$2+\lambda_i$ of $2I+A$, so its determinant is zero exactly when one
factor is zero. On a nonempty finite multiset, every element is at
least $-2$ exactly when its minimum is, and its minimum is $-2$ exactly
when that bound holds and $-2$ occurs. The upper bound is the other
shift in Theorem~\ref{thm:psd}. For the empty graph the first
equivalence is vacuous and both sides of the determinant-zero
equivalence are false, since the empty determinant is $1$.
\end{proof}
\lean{SimpleGraph.posSemidef_two_add_adjMatrix_iff, SimpleGraph.neg_two_mem_spectrumMultiset_iff_det_eq_zero, SimpleGraph.leastEigenvalue_eq_neg_two_iff}
\status{New from M1, M3 and the preceding theorems.}
\used{[CGSS] before Theorem 3.8, Section 4 opening, Theorem 4.13,
Section 5; the standing hypothesis
\texttt{(2 \ensuremath{\bullet} 1 + A).PosSemidef} of the Lean statements.}

\begin{theorem}[Gram matrices with rank control]\label{thm:gram}
A real symmetric matrix $G$ is positive semidefinite iff it is a Gram
matrix. If $G\succeq0$ and $\rk G=\rho$, it has a Gram representation in $\re^\rho$
whose vectors span that space. For any real matrix $B$,
\[
 \ker(B^TB)=\ker B,\qquad \rk(B^TB)=\rk B=\rk(BB^T).
\]
Thus the rank of the Gram matrix of a finite vector family equals the
dimension of its span. In particular $2I+A(X)\succeq0$ iff $X$ has a
representation by vectors of squared norm $2$ with inner product $1$
on edges and $0$ on distinct nonedges. If such vectors lie in dimension
$\ell<n$, then $-2$ occurs, with multiplicity at least $n-\ell$.
For a Gram representation of $2I-A(X)$ in dimension $\ell$, every
eigenvalue is at most $2$ and $\mult_X(2)\ge n-\ell$.
If $n>\ell$, the largest eigenvalue is $2$.
\end{theorem}
\begin{proof}
For a positive semidefinite $G$, choose its orthonormal eigenvectors
and retain just the $\rho$ positive eigenvalues. Define
$B(i,u)=\sqrt{\lambda_i}\Zv_i(u)$ on those rows. Then
$B^TB=\sum_{i:\lambda_i>0}\lambda_i\Zv_i\Zv_i^T=G$.
The rows of $B$ are independent, since they are nonzero multiples of
orthonormal vectors, so $B$ has rank $\rho$ and its columns span
$\re^\rho$. Conversely $\Zx^TB^TB\Zx=\norm{B\Zx}^2\ge0$.
The same identity proves $B^TB\Zx=0\Rightarrow B\Zx=0$ by taking the
inner product with $\Zx$; the reverse implication is direct. Rank-nullity
then gives equal ranks. Applying this to $B^T$, and using equality of
row rank and column rank, gives the other equality. Alternatively,
$\ker B^T=(\im B)^\perp$, since both conditions say that the vector
has zero inner product with every column of $B$; taking dimensions
proves $\rk B^T=\rk B$.
For a family in a finite-dimensional inner product space, take
coordinates in an orthonormal basis of its span. Its coordinate matrix
has image equal to that span and Gram matrix $B^TB$, proving the rank
claim. The entries of $2I+A$ give exactly the stated inner products.
Finally $\rk(2I+A)\le\ell$, whereas
Theorem~\ref{thm:calculus} and Theorem~\ref{thm:trace} give
$\rk(2I+A)=n-\mult_A(-2)$. This proves the bound and occurrence. For $2I-A$, positivity gives
the upper bound by Theorem~\ref{thm:psd}, while
$\mult_A(2)=n-\rk(2I-A)\ge n-\ell$. If $n>\ell$, this count is
positive; Lemma~\ref{lem:extremes} then identifies the maximum.
\end{proof}
\lean{Matrix.IsHermitian.exists_gram_rank, Matrix.rank_gram_eq_finrank_span, SimpleGraph.exists_representation_iff_posSemidef, SimpleGraph.multiplicity_two_ge_of_gram_representation}
\status{Mathlib: M3's Gram existence and rank facts. New: square
factorization, exact-rank representation and graph interface.}
\used{LA1, LA3, [CGSS] Sections 3--4 and 6; representation rank in Theorem 5.8.}

\begin{lemma}[Trace and rank inequality]\label{lem:trace-rank-bound}
Let $G$ be a real positive semidefinite matrix of rank $\rho$. Then
$(\tr G)^2\le\rho\tr(G^2)$, with equality iff all nonzero eigenvalues
are equal; equality includes $G=0$.
\end{lemma}
\begin{proof}
For the $\rho$ nonzero eigenvalues $\eta_i$, expand
\[
 \rho\sum_i\eta_i^2-\Bigl(\sum_i\eta_i\Bigr)^2
 =\sum_{i<j}(\eta_i-\eta_j)^2.
\]
Each term is nonnegative, and the sum is zero exactly when all pairs
are equal. Theorem~\ref{thm:trace} identifies the two sums as traces.
If $\rho=0$, the spectral theorem gives $G=0$ and both sides are zero.
\end{proof}
\lean{Matrix.IsHermitian.trace_sq_le_rank_mul_trace_sq}
\status{New finite-sum inequality, reusable for any positive semidefinite matrix.}
\used{The relative bound for a second Gram matrix.}

\begin{theorem}[Trace bounds and tight frames]\label{thm:trace-bound}
Let $G$ be a real positive semidefinite matrix of rank $\rho$. Then
$(\tr G)^2\le\rho\tr(G^2)$; equality holds iff all its nonzero
eigenvalues are equal, with equality also allowed when $G=0$.
If $n$ unit vectors in $\re^\ell$ have pairwise inner products of
absolute value $\alpha$, where $0\le\alpha\le1$ and
$\ell\alpha^2<1$, then
\[
 n\le\frac{\ell(1-\alpha^2)}{1-\ell\alpha^2}.
\]
If $B$ is a real $\ell\times n$ matrix, $\ell>0$, and
$BB^T=bI$ with $b>0$, then $G=B^TB$ has spectrum
$b^{\ell},0^{n-\ell}$. In particular, with $C=G-2I$,
\[
 (C+2I)(C-(b-2)I)=0.
\]
\end{theorem}
\begin{proof}
Apply Lemma~\ref{lem:trace-rank-bound} to $G$ for the first assertion.
For the Gram matrix of the unit vectors, $\tr G=n$ and
$\tr G^2=\sum_{u,v}G(u,v)^2=n+\alpha^2n(n-1)$.
Its rank is at most $\ell$ by Theorem~\ref{thm:gram}. If $n>0$,
Lemma~\ref{lem:trace-rank-bound} applied to this Gram matrix, divided by $n$, gives
$n(1-\ell\alpha^2)\le\ell(1-\alpha^2)$. Division by its positive
left coefficient gives the bound. If $n=0$, the right side is
nonnegative and the bound holds as well.
For the last assertion, $\rk B=\ell$ by
Theorem~\ref{thm:gram}, so $n\ge\ell$ and $\rk G=\ell$.
Furthermore $G^2=B^T(BB^T)B=bG$. Thus each eigenvalue is $0$ or $b$;
the rank counts precisely $\ell$ occurrences of $b$. Substitution
$G=C+2I$ gives the polynomial identity.
\end{proof}
\lean{Matrix.spectrumMultiset_gram_of_tightFrame}
\status{New wrappers and inequalities from M1, M3 and finite sums.}
\used{EQ1--EQ3; [CGSS] Theorem 3.7, Theorems 4.3--4.4, C6, C24.
For the $E_8$ tight frame $b=30$, the factor is $C-28I$.}

\section{Rayleigh quotients and symmetric Perron--Frobenius}\label{sec:perron}

\begin{theorem}[Rayleigh extrema and equality]\label{thm:rayleigh}
Let $A$ be a real symmetric matrix on a nonempty finite set. For
$\Zx\ne0$, put $R_A(\Zx)=\Zx^TA\Zx/\norm{\Zx}^2$. Then
\[
 \lambda_{\min}(A)\le R_A(\Zx)\le\lambda_{\max}(A).
\]
In any orthonormal eigenbasis, writing $\Zx=\sum_i b_i\Zv_i$ gives
\[
 R_A(\Zx)=\frac{\sum_i\lambda_i b_i^2}{\sum_i b_i^2},\qquad
 \sum_i b_i^2=\norm{\Zx}^2>0,
\]
\[
 \begin{aligned}
 \lambda_{\max}\norm{\Zx}^2-\Zx^TA\Zx
   &=\sum_i(\lambda_{\max}-\lambda_i)b_i^2,\\
 \Zx^TA\Zx-\lambda_{\min}\norm{\Zx}^2
   &=\sum_i(\lambda_i-\lambda_{\min})b_i^2.
 \end{aligned}
\]
Both bounds are attained, including on the unit sphere. Equality on
the right holds iff $A\Zx=\lambda_{\max}(A)\Zx$, and equality on
the left holds iff $A\Zx=\lambda_{\min}(A)\Zx$. Thus the extremes
are the supremum and infimum of the Rayleigh quotient on nonzero vectors.
Eigenvectors of a symmetric matrix for distinct eigenvalues are orthogonal.
\end{theorem}
\begin{proof}
In an orthonormal eigenbasis,
$R_A(\Zx)=\sum_i\lambda_i b_i^2/\sum_i b_i^2$; its denominator is
positive. This is a weighted average of the eigenvalues. More
explicitly,
$\lambda_{\max}\norm{\Zx}^2-\Zx^TA\Zx
=\sum_i(\lambda_{\max}-\lambda_i)b_i^2$.
Lemma~\ref{lem:extremes} makes the terms nonnegative, and vanishing is equivalent to $b_i=0$
whenever $\lambda_i<\lambda_{\max}$, which is the stated eigenvector
condition. For the minimum, subtract
$\lambda_{\min}\norm{\Zx}^2$ from the quadratic form; the terms are
$(\lambda_i-\lambda_{\min})b_i^2$, and the identical zero-term
argument gives the left equality condition. A unit basis vector for
each extreme attains it, so each is the claimed supremum or infimum.
If $A\Zx=s\Zx$ and $A\Zy=t\Zy$, symmetry gives
$s\ip{\Zx}{\Zy}=\ip{A\Zx}{\Zy}=\ip{\Zx}{A\Zy}
=t\ip{\Zx}{\Zy}$. For $s\ne t$, cancellation gives orthogonality.
\end{proof}
\lean{Matrix.IsHermitian.rayleigh_le_largestEigenvalue, Matrix.IsHermitian.rayleigh_eq_largestEigenvalue_iff, Matrix.IsHermitian.leastEigenvalue_le_rayleigh}
\status{Mathlib: M5 at the linear-map level. New: these matrix
wrappers and an explicit equality criterion from M1.}
\used{SP2, SM1, [CGSS] Theorem 3.8 and Lemma 5.2.}

\begin{lemma}[Rayleigh subspace witnesses]\label{lem:rayleigh-subspaces}
Let $A$ be real symmetric of order $n>0$, with the sorted eigenbasis of
Theorem~\ref{thm:multiplicity} and eigenvalues $\nu_0\ge\cdots\ge\nu_{n-1}$.
For $0\le i<n$, the head span of directions $0,\ldots,i$ has dimension
$i+1$ and quotient at least $\nu_i$ on nonzero vectors. The tail span
of directions $i,\ldots,n-1$ has dimension $n-i$ and quotient at most
$\nu_i$. Every subspace of dimension $i+1$ contains a nonzero vector
with quotient at most $\nu_i$; every subspace of dimension $n-i$
contains one with quotient at least $\nu_i$.
\end{lemma}
\begin{proof}
Independence of the sorted basis gives the two span dimensions.
The coordinate quotient in Theorem~\ref{thm:rayleigh} is a weighted
average of entries bounded below by $\nu_i$ on the head, and bounded
above by $\nu_i$ on the tail. For subspaces $P,Q$ of $\re^n$,
$\dim(P\cap Q)=\dim P+\dim Q-\dim(P+Q)$ and $\dim(P+Q)\le n$.
Intersect an $(i+1)$-dimensional subspace with the tail: the dimension
sum is $n+1$, so their intersection has a nonzero vector. The head
intersects any $(n-i)$-dimensional subspace by the same dimension
formula. These vectors give the two witnesses.
\end{proof}
\lean{Matrix.IsHermitian.rayleigh_head_tail, Matrix.IsHermitian.exists_rayleigh_le_of_finrank, Matrix.IsHermitian.exists_le_rayleigh_of_finrank}
\status{New witness form of Courant--Fischer from M1 and dimensions.}
\used{Both applications, to the larger and smaller matrices, in interlacing.}

\begin{theorem}[Cauchy interlacing]\label{thm:interlacing}
Let $A$ be real symmetric of order $n$, and let $B$ be its principal
submatrix on $m$ distinct indices, where $1\le m\le n$. Write
$\lambda_0\ge\cdots\ge\lambda_{n-1}$ and
$\beta_0\ge\cdots\ge\beta_{m-1}$ for the decreasing eigenvalue lists
$\nu_i$ of $A$ and of $B$ in Theorem~\ref{thm:multiplicity}; they list
the two spectra with repetitions. Then
\[
 \lambda_i\ge\beta_i\ge\lambda_{i+n-m}\qquad(0\le i<m).
\]
In particular the largest eigenvalue cannot increase, and the least
eigenvalue cannot decrease, on taking a nonempty principal submatrix.
\end{theorem}
\begin{proof}
Use the sorted bases of Theorem~\ref{thm:multiplicity} for $A$ and $B$;
their eigenvalues are the two lists. Extension by zero embeds $\re^m$ into $\re^n$, preserving dimension,
norm and quadratic form. By Lemma~\ref{lem:rayleigh-subspaces} for
$B$, its extended head of dimension $i+1$ has quotient at least
$\beta_i$. Apply that lemma to $A$ and this subspace to obtain a
nonzero vector with quotient at most $\lambda_i$. Thus
$\beta_i\le\lambda_i$.
The extended tail for $B$ has dimension $m-i$ and quotient at most
$\beta_i$. Apply the second witness assertion for $A$ at index
$i+n-m$, since $m-i=n-(i+n-m)$. It gives a quotient at least
$\lambda_{i+n-m}$, proving the other inequality. At $i=0$ and
$i=m-1$, Lemma~\ref{lem:extremes} identifies the endpoint bounds.
\end{proof}
\lean{Matrix.IsHermitian.eigenvalues_principal_submatrix_le, Matrix.IsHermitian.le_eigenvalues_principal_submatrix}
\status{New proof from M1 and finite-dimensional subspace arithmetic.}
\used{Spectral bounds for represented subsets in [CGSS] Section 4.}

\begin{definition}\label{def:support}
For a symmetric entrywise nonnegative matrix $A$, its \emph{positive
support graph} has the same indices as vertices and an edge $uv$ when
$u\ne v$ and $A(u,v)>0$. Connectedness means that the vertex set is
nonempty and any two vertices are joined by a path. In particular a
single vertex with no edges is connected.
\lean{Matrix.positiveSupportGraph}
\status{New graph definition; it uses \code{SimpleGraph.Connected}.
It can later be compared with the quiver definition in M5.}
\end{definition}

\begin{theorem}[Symmetric Perron--Frobenius]\label{thm:perron}
Let $A$ be a real symmetric entrywise nonnegative matrix on a nonempty
finite set. Put $\rho=\lambda_{\max}(A)$. Then $\rho\ge0$, there is
a nonzero entrywise nonnegative vector $\Zx$ with $A\Zx=\rho\Zx$,
and $|\mu|\le\rho$ for every eigenvalue $\mu$.
If the positive support graph is connected, this vector is strictly
positive in every coordinate, $\mult_A(\rho)=1$, and every
nonzero nonnegative $\rho$-eigenvector is strictly positive.
\end{theorem}
\begin{proof}
For any coordinate unit vector the Rayleigh quotient is a diagonal
entry of $A$, thus nonnegative; Theorem~\ref{thm:rayleigh} gives
$\rho\ge0$. Choose a unit $\rho$-eigenvector $\Zy$ and let
$\Zx(u)=|\Zy(u)|$. Its norm is one and, term by term,
\[
 \Zy^TA\Zy=\sum_{u,v}A(u,v)\Zy(u)\Zy(v)
 \le\sum_{u,v}A(u,v)|\Zy(u)||\Zy(v)|=\Zx^TA\Zx.
\]
The first expression is $\rho$ and the last is at most $\rho$ by
Rayleigh. Thus equality holds and $A\Zx=\rho\Zx$ by its equality
criterion. For an arbitrary $\mu$-eigenvector $\Zw\ne0$, the triangle
inequality and Rayleigh give
$|\mu|\norm{\Zw}^2=|\Zw^TA\Zw|
\le|\Zw|^TA|\Zw|\le\rho\norm{\Zw}^2$.
Division proves the absolute-value bound.

Now assume connected support. If a nonnegative eigenvector $\Zx$ has
$\Zx(u)=0$, its eigenvector equation at $u$ says
$\sum_v A(u,v)\Zx(v)=0$. Each summand is nonnegative, so $\Zx(v)=0$
at each positive-support neighbour. Induction along a path propagates
zero to every coordinate, a contradiction to $\Zx\ne0$. Thus such
an eigenvector is positive.
Fix a positive $\rho$-eigenvector $\Zx$, and let $\Zy$ be any real
$\rho$-eigenvector, allowing zero. The finite nonempty set of ratios
$\Zy(u)/\Zx(u)$ has a minimum $b$, attained at some $u$. The vector
$\Zy-b\Zx$ is nonnegative, is a $\rho$-eigenvector, and has zero
coordinate at $u$. The propagation just proved forces it to be zero.
Every vector in the eigenspace is thus a scalar multiple of
$\Zx$. Its dimension, and thus its multiplicity by
Theorem~\ref{thm:multiplicity}, is one.
\end{proof}
\lean{Matrix.IsHermitian.exists_nonneg_eigenvector_largest, Matrix.IsHermitian.exists_pos_eigenvector_largest, Matrix.IsHermitian.multiplicity_largest_eq_one, Matrix.IsHermitian.abs_eigenvalue_le_largest}
\status{New from the spectral theorem and Rayleigh equality. No
Perron--Frobenius theorem was found in Mathlib.}
\used{SM1, Theorem 3.8, C17, and weighted symmetric matrices in Remark 3.11.}

\begin{lemma}[A positive supervector bounds each eigenvalue]\label{lem:supervector}
Let $A$ be an entrywise nonnegative real matrix on a nonempty finite set,
$\Zx>0$ coordinatewise, and $A\Zx\le b\Zx$. If
$A\Zy=\theta\Zy$ with $\Zy\ne0$, then $\theta\le b$.
Symmetry is not required.
\end{lemma}
\begin{proof}
Choose $u$ maximizing $|\Zy(u)|/\Zx(u)$ and call this maximum $c$.
A nonzero coordinate of $\Zy$ implies $c>0$; thus
$|\Zy(u)|=c\Zx(u)>0$. Every $v$ satisfies $|\Zy(v)|\le c\Zx(v)$.
The triangle inequality and nonnegativity give
\[
 \theta|\Zy(u)|\le |(A\Zy)(u)|
 \le\sum_v A(u,v)|\Zy(v)|\le c(A\Zx)(u)
 \le cb\Zx(u)=b|\Zy(u)|.
\]
Divide by $|\Zy(u)|>0$.
\end{proof}
\lean{Matrix.eigenvalue_le_of_pos_supervector}
\status{Implemented with this maximum-ratio proof in
\code{LineGraphs/Spectral/PositiveVector.lean}.}
\used{Positive-vector certificates without a Perron--Frobenius dependency.}

\begin{corollary}[Positive-vector certificates]\label{cor:positive-certificate}
Let $A$ be a real symmetric nonnegative matrix on a nonempty finite
set and let every coordinate of $\Zx$ be positive. If
$A\Zx=b\Zx$, then $b=\lambda_{\max}(A)$, whether or not its support
is connected. More generally, $A\Zx\le b\Zx$ coordinatewise implies
$\lambda_{\max}(A)\le b$.
\end{corollary}
\begin{proof}
Apply Lemma~\ref{lem:supervector} to a largest-eigenvalue basis vector
from Theorem~\ref{thm:multiplicity} and Lemma~\ref{lem:extremes}.
This gives the upper bound. If $A\Zx=b\Zx$, positivity and nonemptiness
make $\Zx\ne0$, so $b$ belongs to the spectrum by
Theorem~\ref{thm:multiplicity}. Lemma~\ref{lem:extremes} gives the
reverse inequality. This is the proof in
\code{LineGraphs/Spectral/PositiveVector.lean}.
\end{proof}
\lean{Matrix.IsHermitian.largestEigenvalue_eq_of_pos_eigenvector, Matrix.IsHermitian.largestEigenvalue_le_of_pos_supervector}
\status{Implemented maximum-ratio consequence in \code{LineGraphs/Spectral/PositiveVector.lean}.}
\used{SM2 and C17: an integer vector identity and positive coordinates suffice.}

\begin{theorem}[Strict monotonicity]\label{thm:strict}
Let $A$ be real symmetric and nonnegative with connected positive
support. Let $B$ be real symmetric and nonnegative on a nonempty
finite set $W$, and let $\phi:W\hookrightarrow U$ be injective with
$B(u,v)\le A(\phi(u),\phi(v))$ for all $u,v$.
Then $\lambda_{\max}(B)\le\lambda_{\max}(A)$.
The inequality is strict if $\phi$ is not surjective, or if one of
these entry inequalities is strict. In particular every nonempty
proper subgraph of a connected finite graph has strictly smaller
largest eigenvalue. A subgraph here may omit edges as well as vertices.
\end{theorem}
\begin{proof}
Take a nonnegative unit top eigenvector of $B$, and extend it by zero
using $\phi$ to a unit vector $\Zz$ on $U$. Nonnegativity gives
$\lambda_{\max}(B)\le\Zz^TA\Zz\le\lambda_{\max}(A)$.
If the extremes are equal, Rayleigh equality makes $\Zz$ a top
eigenvector of $A$; it is positive by connectedness and
Theorem~\ref{thm:perron}. Thus no coordinate was omitted.
Every coordinate of the original eigenvector is now positive.
The difference of the two quadratic forms is
$\sum_{u,v}(A(\phi(u),\phi(v))-B(u,v))\Zz(\phi(u))\Zz(\phi(v))$.
All terms are nonnegative; a strict entry inequality makes one term
positive. This contradicts equality. For graphs, a vertex injection
which preserves edges gives precisely these inequalities between
adjacency matrices, and a proper subgraph has one of the two strictness
conditions.
\end{proof}
\lean{Matrix.IsHermitian.largestEigenvalue_lt_of_lt_submatrix, SimpleGraph.largestEigenvalue_lt_of_proper_subgraph}
\status{New; the non-strict principal-submatrix case also follows from interlacing.}
\used{SM1, SM3 and [CGSS] Theorem 3.8. The smaller graph must be nonempty.}

\section{Walks, Laplacians, and regular graphs}\label{sec:regular}

\begin{theorem}[Walks and moments]\label{thm:walks}
Let $X$ be a finite simple graph with $n$ vertices, $m$ edges, and
adjacency matrix $A$. For $j\ge0$, $(A^j)(u,v)$ counts walks of
length $j$ from $u$ to $v$, and $\tr A^j$ counts closed walks with
a distinguished starting vertex. In particular
\[
 \sum_{\mu\in\Spec(X)}1=n,\quad
 \sum_{\mu\in\Spec(X)}\mu=0,\quad
 \sum_{\mu\in\Spec(X)}\mu^2=2m=\sum_u\deg(u),\quad
 \sum_{\mu\in\Spec(X)}\mu^3=6t,
\]
where $t$ is the number of triangles and all sums count multiplicity.
At each vertex $u$,
$(A^3)(u,u)=2\,\#\{\text{triangles containing }u\}$.
The total number of walks of length $j$ is $\Zj^TA^j\Zj$; it need
not equal $\tr A^j$.
\end{theorem}
\begin{proof}
For $j=0$ the identity matrix counts the unique length-zero walk from
a vertex to itself and no walks between distinct vertices. A walk of
length $j+1$ from $u$ to $v$ is specified uniquely by its first
neighbour $w$ of $u$ and a length-$j$ walk from $w$ to $v$.
Thus its count is $\sum_wA(u,w)(A^j)(w,v)$, proving the assertion
by induction. Sum over $u=v$ to obtain the closed-walk count and
over all $u,v$ for the total-walk count. There are no closed walks
of length one. A closed walk of length two chooses a vertex and an
incident edge; each edge gives two such choices. In length three,
no two consecutive vertices can coincide and the first three vertices
are distinct, because otherwise one of the three edges is a loop.
Thus the walk traverses a triangle. Each triangle gives three choices
of starting vertex and two directions, exactly six walks. The
same argument with the starting vertex fixed gives two walks per
triangle through that vertex. The trace-power formula of
Theorem~\ref{thm:trace} gives all four
spectral identities. The handshake equality counts the same incident
vertex-edge pairs by vertices or by edges.
\end{proof}
\lean{SimpleGraph.trace_pow_adjMatrix_eq_closedWalks, SimpleGraph.sum_sq_spectrumMultiset_eq_twice_card_edges, SimpleGraph.trace_cube_adjMatrix_eq_six_mul_triangles, SimpleGraph.adjMatrix_cube_apply_self_eq_twice_triangles}
\status{Mathlib: the entry count and first two diagonal calculations in M6.
New: the second moment and total-walk identities, multiset wrappers
and the triangle trace interface.}
\used{SP2, [CGSS] Theorem 3.7, Theorem 5.8, C13--C15.}

\begin{theorem}[Laplacian and signless Laplacian kernels]\label{thm:laplacian}
Let $X$ be a finite simple graph, let $D$ be the diagonal degree matrix,
and put $L=D-A$ and $K=D+A$. Then
\[
 \Zx^TL\Zx=\sum_{\{u,v\}\in E(X)}(\Zx(u)-\Zx(v))^2,\qquad
 \Zx^TK\Zx=\sum_{\{u,v\}\in E(X)}(\Zx(u)+\Zx(v))^2.
\]
Both matrices are positive semidefinite. Their kernels consist,
respectively, of functions constant across each edge and functions
whose values sum to zero across each edge. Thus
\[
 \dim\ker L=\#\{\text{components of }X\},\qquad
 \dim\ker K=\#\{\text{bipartite components of }X\}.
\]
An isolated vertex counts as a bipartite component.
\end{theorem}
\begin{proof}
In each sum of squares, the square $\Zx(u)^2$ occurs once for each
edge at $u$, so its coefficient is $\deg(u)$. The cross terms are
$\mp2\sum_{\{u,v\}}\Zx(u)\Zx(v)$, while
$\Zx^TA\Zx=2\sum_{\{u,v\}}\Zx(u)\Zx(v)$, proving the formulas.
They imply positivity. By Theorem~\ref{thm:psd}, membership in a
kernel is equivalent to vanishing of the corresponding quadratic
form, which is equivalent to each square vanishing.
For $L$, constancy propagates along paths; arbitrary constants on
the components give all kernel functions. The characteristic
functions of the components are a basis: evaluate a linear relation
at a vertex of each component to prove independence, and choose the
constant values to prove spanning.

For $K$, fix a connected component and a vertex $u$ in it. The edge
relation determines the value at the end of any path of length $j$
to be $(-1)^j\Zx(u)$. In particular the space of such functions has
dimension at most one, since a function zero at $u$ is zero everywhere
on the component. If a nonzero function exists, all its absolute
values equal $|\Zx(u)|>0$. Partition the vertices by positive and
negative value. The edge relation makes every edge cross the
partition, so the component is bipartite. Conversely a bipartition
gives the function $1$ on one part and $-1$ on the other, a nonzero
kernel function; for an isolated vertex take the first part to be
the vertex. This proves dimension one exactly on the bipartite
components and zero on the others. Extend these signed functions by
zero to the remaining components. Disjoint supports prove their
independence, and the connected-component argument proves spanning.
\end{proof}
\lean{SimpleGraph.signlessLapMatrix, SimpleGraph.signlessLapMatrix_mulVec_eq_zero_iff, SimpleGraph.finrank_ker_signlessLapMatrix_eq_card_bipartite_components}
\status{Mathlib: the Laplacian half in M6. New: the signless matrix
and its component-count theorem.}
\used{SP1, SP3, SP5, [CGSS] Lemmas 5.1, 5.3, 5.5--5.6.}

\begin{theorem}[SP1: regular graphs]\label{thm:sp1}
Let $X$ be a nonempty finite $k$-regular graph. Its largest eigenvalue
is $k$, every eigenvalue lies in $[-k,k]$, and
\[
 \mult_X(k)=\#\{\text{connected components of }X\}.
\]
\end{theorem}
\begin{proof}
The row sum of $A$ is $k$, so $A\Zj=k\Zj$ and $\Zj\ne0$.
Here $L=kI-A$ and $K=kI+A$ are positive semidefinite by
Theorem~\ref{thm:laplacian}. The shift criterion in
Theorem~\ref{thm:psd} gives $-k\le\mu\le k$ for every eigenvalue;
since $k$ occurs it is the maximum. The equation $A\Zx=k\Zx$ is
equivalent to $L\Zx=0$. Multiplicity equals its kernel dimension
by Theorem~\ref{thm:multiplicity}, and
Theorem~\ref{thm:laplacian} supplies the component count.
\end{proof}
\lean{SimpleGraph.IsRegularOfDegree.largestEigenvalue_eq, SimpleGraph.IsRegularOfDegree.multiplicity_eq_card_components}
\status{New graph bridge from M6 and the multiplicity theorem.}
\used{[CGSS] Lemma 5.1; regularity and connectedness deductions throughout Section 5.}

\begin{theorem}[SP2: regularity detected by the spectrum]\label{thm:sp2}
For a nonempty finite graph $X$ on $n$ vertices,
\[
 X\text{ is regular}\ \Longleftrightarrow\
 \sum_{\mu\in\Spec(X)}\mu^2=n\lambda_{\max}(X).
\]
\end{theorem}
\begin{proof}
We have $\norm{\Zj}^2=n$ and
$\Zj^TA\Zj=\sum_u\deg(u)=\tr A^2$ by
Theorem~\ref{thm:walks}. Thus the displayed identity is precisely
$R_A(\Zj)=\lambda_{\max}(A)$, with $n>0$. Rayleigh equality says
this is equivalent to $A\Zj=\lambda_{\max}\Zj$, or
$\deg(u)=\lambda_{\max}$ at every vertex. Choose a vertex $u$;
its natural-number degree $k$ then equals this real value, and all
degrees equal $k$. This proves regularity and the reverse implication
as well.
\end{proof}
\lean{SimpleGraph.isRegular_iff_sum_sq_eigenvalues}
\status{New from M1, M5, M6 and the second-moment identity of
Theorem~\ref{thm:walks}.}
\used{[CGSS] Lemma 5.2 and the implicit connectedness step in Theorem 5.8.}

\begin{corollary}[Cospectral mates of regular graphs]\label{cor:regular-mates}
Every graph cospectral with a nonempty finite $k$-regular graph is
$k$-regular, has the same number of components, and is connected if
that graph is connected.
\end{corollary}
\begin{proof}
Proposition~\ref{prop:cospectral} gives the same nonzero order, second
moment and largest eigenvalue. Apply Theorem~\ref{thm:sp2} to the mate.
Its degree is the common largest eigenvalue $k$ by
Theorem~\ref{thm:sp1}. The common multiplicity of $k$ gives the same
component count by that theorem. One component is equivalent to
connectedness for a nonempty graph.
\end{proof}
\lean{SimpleGraph.Cospectral.isRegularOfDegree, SimpleGraph.Cospectral.connected_of_regular}
\status{New corollary applying SP2 to each graph separately.}
\used{[CGSS] Lemma 5.2 and cospectral line-graph arguments.}

\begin{theorem}[SP3: bipartiteness]\label{thm:sp3}
Let $X$ be a connected finite $k$-regular graph. It is bipartite iff
$-k\in\Spec(X)$; when this occurs its multiplicity is one.
More generally, for every finite bipartite graph $Y$,
\[
 \mult_Y(\mu)=\mult_Y(-\mu)\qquad(\mu\in\re).
\]
\end{theorem}
\begin{proof}
The $-k$-eigenspace is $\ker(kI+A)$. By
Theorem~\ref{thm:laplacian} its dimension for a connected graph is
one if the graph is bipartite and zero otherwise. Apply
Theorem~\ref{thm:multiplicity}. This includes the connected graph
with one vertex and $k=0$: both assertions refer to its single
zero eigenvalue. For the second assertion, let $T$ be the diagonal
matrix of signs $1$ and $-1$ on the two parts. Then $T^2=I$ and
$TAT=-A$, since every edge crosses parts and zero entries remain zero.
If $A\Zx=\mu\Zx$, then $A(T\Zx)=-\mu T\Zx$. Multiplication by
$T$ is its own inverse, so it restricts to a linear equivalence of
the two eigenspaces. Their dimensions and multiplicities agree.
\end{proof}
\lean{SimpleGraph.IsRegularOfDegree.isBipartite_iff_neg_eigenvalue, SimpleGraph.IsBipartite.multiplicity_neg_eq}
\status{New signless-kernel bridge.}
\used{[CGSS] Lemma 5.3 and the nonsingularity step of Lemma 5.6.}

\begin{theorem}[The all-ones direction]\label{thm:ones}
Let $A$ be real symmetric of nonzero order $n$ and satisfy $A\Zj=k\Zj$.
For every $\Zx$, $J\Zx=(\sum_u\Zx(u))\Zj$; thus $J\Zx=0$ on
$\Zj^\perp$. There is an orthogonal decomposition
$\re^U=\operatorname{span}(\Zj)\oplus\Zj^\perp$, of dimensions $1,n-1$.
The orthogonal complement $\Zj^\perp$ is invariant under $A$, so one
can choose an orthonormal eigenbasis starting with $\Zj/\sqrt n$.
Write the remaining eigenvalues, with multiplicity, as $\mu_1,\ldots,
\mu_{n-1}$. For real $b,s,t$, the eigenvalues of $bI+sA+tJ$ are
\[
 b+sk+tn,\qquad b+s\mu_i\quad(1\le i<n).
\]
For every finite simple graph $X$, $A(\comp X)=J-I-A(X)$.
Thus for a nonempty $k$-regular graph,
\[
 \begin{aligned}
 \Spec(\comp X)&=\ms{n-1-k,-1-\mu_1,\ldots,-1-\mu_{n-1}},\\
 \Spec(S(X))&=\ms{n-1-2k,-1-2\mu_1,\ldots,-1-2\mu_{n-1}}.
 \end{aligned}
\]
\end{theorem}
\begin{proof}
If $\ip{\Zx}{\Zj}=0$, symmetry gives
$\ip{A\Zx}{\Zj}=\ip{\Zx}{A\Zj}=0$, proving invariance.
The restriction of $A$ to this finite-dimensional inner product space
is symmetric; apply the spectral theorem there and adjoin the
normalized all-ones vector. The matrix $J$ acts by $n$ on that
vector and by zero on $\Zj^\perp$, since
$J\Zx=(\sum_u\Zx(u))\Zj$. Applying the three summands to this
basis proves the eigenvalue list. Mathlib's
\code{SimpleGraph.one_add_adjMatrix_add_compl_adjMatrix_eq_of_one} and
\code{SimpleGraph.compl_adjMatrix_eq_adjMatrix_compl} say that
$I+A(X)+A(\comp X)=J$. Substitute $A(\comp X)=J-I-A$
and $S=J-I-2A$ for the graph claims.
\end{proof}
\lean{Matrix.IsHermitian.spectrumMultiset_smul_add_smul_one_add_smul_allOnes, SimpleGraph.IsRegularOfDegree.spectrumMultiset_compl, SimpleGraph.IsRegularOfDegree.spectrumMultiset_seidelMatrix}
\status{Mathlib: the complement identity in M6. New wrappers. An
eigenvector for an eigenvalue different from $k$ is orthogonal to $\Zj$;
here extra copies of $k$ in $\Zj^\perp$ are retained.}
\used{SR2, [CGSS] Theorem 4.4, Corollary 5.10, Section 6.}

\begin{corollary}[Removing the constant direction]\label{cor:remove-ones}
Let $G\succeq0$ be real symmetric of order $n>0$, with $G\Zj=b\Zj$
and $b>0$. Then $G-(b/n)J\succeq0$ and its rank is $\rk G-1$.
\end{corollary}
\begin{proof}
Apply Theorem~\ref{thm:ones} to $G$ with coefficients $0,1,-b/n$.
Its eigenvalues are zero on the constant direction and the unchanged
nonnegative eigenvalues of $G$ on its orthogonal complement.
Theorems~\ref{thm:psd} and \ref{thm:trace} give positivity and rank:
exactly one nonzero eigenvalue, namely $b$, was removed.
\end{proof}
\lean{Matrix.PosSemidef.rank_sub_allOnes}
\status{New application of the all-ones spectrum to a second matrix.}
\used{[CGSS] Theorem 4.4 and centred Gram representations.}

\begin{corollary}[A scalar matrix plus the all-ones matrix]\label{cor:scalar-ones}
For real $b,t$ and order $n>0$,
\[
 \begin{aligned}
 \Spec(bI+tJ)&=\ms{(b+tn)^1,b^{n-1}},\\
 \chi_{bI+tJ}(z)&=(z-b-tn)(z-b)^{n-1},\\
 \det(bI+tJ)&=(b+tn)b^{n-1}.
 \end{aligned}
\]
The multiset convention combines the two values when $t=0$.
\end{corollary}
\begin{proof}
The all-ones vector is nonzero and $J\Zj=n\Zj$. Every vector with
coordinate sum zero is killed by $J$. These two orthogonal spaces
have dimensions one and $n-1$, as in Theorem~\ref{thm:ones}.
Thus $bI+tJ$ acts by $b+tn$ on the first and by $b$ on the second.
Theorem~\ref{thm:multiplicity} gives the spectrum and characteristic
polynomial, and Theorem~\ref{thm:trace} gives the determinant.
This proof includes $b=0$ and $n=1$, with a zeroth power interpreted
as one. For order zero the characteristic polynomial and determinant
are both one; the displayed formula is deliberately restricted to
$n>0$.
\end{proof}
\lean{Matrix.spectrumMultiset_smul_one_add_smul_allOnes, Matrix.charpoly_smul_one_add_smul_allOnes, Matrix.det_smul_one_add_smul_allOnes}
\status{Mathlib: M8 has the rank-one determinant lemma.
New: these explicit real spectrum and determinant wrappers.}
\used{LA6, DS3, DS5 and the determinant calculation in C19.}

\section{Rectangular products and line graphs}\label{sec:line}

\begin{theorem}[Rectangular characteristic polynomials]\label{thm:mn}
Let $M$ be $p\times q$ and $N$ be $q\times p$ over a commutative ring.
Then
\[
 t^q\chi_{MN}(t)=t^p\chi_{NM}(t).
\]
Over $\re$, the nonzero roots have equal algebraic multiplicities and
\[
 q+\mult_{MN}(0)=p+\mult_{NM}(0).
\]
For $\mu\ne0$, multiplication by $N$ is a linear equivalence
$V_{MN}(\mu)\simeq V_{NM}(\mu)$, with inverse $\mu^{-1}M$.
If $B$ is a real $p\times q$ matrix of rank $\rho$, both $BB^T$ and
$B^TB$ are positive semidefinite; their positive spectra agree with
multiplicity, and their zero multiplicities are $p-\rho$ and $q-\rho$.
\end{theorem}
\begin{proof}
Work over the polynomial ring of the original coefficient ring and
form
\[
 G=\pmat{tI_p&M\\N&I_q},\qquad
 T=\pmat{-I_p&0\\N&-tI_q}.
\]
Block multiplication gives
\[
 GT=\pmat{MN-tI_p&-tM\\0&-tI_q},\qquad
 TG=\pmat{-tI_p&-M\\0&NM-tI_q}.
\]
An upper block triangular determinant is the product of its diagonal
block determinants: in the permutation formula, a nonzero term must
map every lower-block row to a lower-block column, and then every
upper-block row to an upper-block column. Thus
$\det(GT)=(-1)^{p+q}t^q\chi_{MN}(t)$ and
$\det(TG)=(-1)^{p+q}t^p\chi_{NM}(t)$.
Multiplicativity and commutativity give $\det(GT)=\det(TG)$.
Cancel the unit $(-1)^{p+q}$ to obtain the polynomial identity.
For real coefficients, root multiplicities add in products. At
$\mu\ne0$ the power of $t$ has multiplicity zero; at zero it has
multiplicity its exponent. This gives both multiplicity assertions,
without any splitting hypothesis on $MN$ or $NM$.

For $MN\Zx=\mu\Zx$, we have
$NM(N\Zx)=N(MN\Zx)=\mu N\Zx$ and
$\mu^{-1}M(N\Zx)=\Zx$. The reverse composite is the identity on
$V_{NM}(\mu)$ by the same eigenvector equation for $NM$.
These are linear inverse maps.
For the Gram products, positivity and ranks follow from
Theorem~\ref{thm:gram}. Symmetric multiplicity and rank-nullity give
the two zero counts. Their remaining eigenvalues are positive by
Theorem~\ref{thm:psd}, and their multiplicities agree by the
characteristic-polynomial identity. All formulas include $p=0$ or
$q=0$; then the asserted nonzero spectra are empty.
\end{proof}
\lean{Matrix.eigenvalueMultiplicity_mul_comm_of_ne_zero, Matrix.spectrumMultiset_gram_products}
\status{Mathlib: \code{Matrix.charpoly_mul_comm'} and the rank facts
(M3--M4). New: multiset and eigenspace-transfer interfaces.}
\used{SP5, [CGSS] Lemmas 5.5--5.7, Theorem 5.8, and Section 6.}

\begin{theorem}[Incidence identities]\label{thm:incidence}
Let $X$ be a finite simple graph with $n$ vertices and $m$ edges.
Define $M(u,e)=1$ if the vertex $u$ is an endpoint of the edge $e$,
and $0$ otherwise, with columns indexed by the actual edges. Then
\[
 MM^T=D+A(X),\qquad M^TM=2I+A(L(X)).
\]
If $b$ is the number of bipartite connected components of $X$, then
\[
 \rk M=n-b,\qquad
 \mult_{L(X)}(-2)=m-n+b.
\]
The right side is nonnegative. Every line graph has least eigenvalue
at least $-2$ whenever it is nonempty.
\end{theorem}
\begin{proof}
The entry $(MM^T)(u,v)$ counts edges containing both vertices. For
$u=v$ it is $\deg(u)$; for distinct $u,v$ it is one on an edge and
zero otherwise. The entry $(M^TM)(e,f)$ counts common endpoints.
For $e=f$ it is two, since there are no loops. Two distinct edges in
a simple graph have at most one common endpoint, with one exactly
when they are adjacent in the line graph. This proves both identities.
By Theorem~\ref{thm:laplacian}, $MM^T$ has kernel dimension $b$.
Its rank is $n-b$, equal to $\rk M$ by Theorem~\ref{thm:gram}.
The other Gram product has rank $n-b$ and size $m$, so its zero
multiplicity is $m-(n-b)$. Its eigenvalues are those of the line
graph shifted by $2$, proving the formula and its nonnegativity.
Finally $M^TM\succeq0$ gives the claimed lower bound by
Corollary~\ref{cor:minus-two}.
\end{proof}
\lean{SimpleGraph.edgeIncMatrix, SimpleGraph.edgeIncMatrix_mul_transpose, SimpleGraph.transpose_edgeIncMatrix_mul, SimpleGraph.multiplicity_neg_two_lineGraph}
\status{Mathlib: \code{SimpleGraph.incMatrix_mul_transpose} on all
unordered pairs (M6). New: restriction to edges, the line-graph
identity, and multiplicity.}
\used{SP5, [CGSS] Section 4 opening and Lemmas 5.5--5.8.}

\begin{lemma}[Connectivity of a line graph]\label{lem:line-connected}
Let $X$ be a finite graph with at least one edge and no isolated
vertices. Then $X$ is connected iff $L(X)$ is connected. For any edge
$\{u,v\}$, its degree in the line graph is $\deg_X(u)+\deg_X(v)-2$.
\end{lemma}
\begin{proof}
If $X$ is connected, take any two edges. A path between an endpoint
of the first and an endpoint of the second, together with the first
and last edges, gives a sequence of edges in which consecutive
distinct edges share an endpoint. Delete consecutive repetitions.
This is a walk in $L(X)$ joining the chosen edges, so $L(X)$ is
connected and nonempty. Conversely choose incident edges at any two
vertices of $X$, using the absence of isolates. A walk joining these
edges in $L(X)$ is a chain of edges of $X$ with intersecting consecutive
endpoint sets. Both endpoints of one edge are reachable from either
endpoint, and transitivity of reachability propagates through the
chain. Thus the two vertices are reachable and $X$ is connected.
For the degree, the edges other than $\{u,v\}$ at $u$ and at $v$
give $\deg(u)-1$ and $\deg(v)-1$ neighbours. These two sets are
disjoint: a different edge with both endpoints $u,v$ would equal the
given edge. Their union is exactly its neighbourhood in $L(X)$.
\end{proof}
\lean{SimpleGraph.lineGraph_connected_iff_of_no_isolated, SimpleGraph.degree_lineGraph}
\status{New elementary interfaces over \code{SimpleGraph.lineGraph}.}
\used{The hypotheses and degree counts of SP4 and [CGSS] Lemmas 5.5--5.6.}

\begin{theorem}[SP4: regular line graphs]\label{thm:sp4}
Let $X$ have at least one edge and no isolated vertices. If $L(X)$
is connected and $h$-regular, then $X$ is either regular or bipartite
with constant degree on each part. In the latter case, writing the
part sizes as $p,q$ and the respective degrees as $k_1,k_2$, we have
$k_1+k_2=h+2$ and $pk_1=qk_2=|E(X)|$.
\end{theorem}
\begin{proof}
The preceding lemma gives connectedness of $X$ and
$\deg(u)+\deg(v)=h+2$ on every edge. Fix a vertex of degree $k_1$
and put $k_2=h+2-k_1$; it has a neighbour, so $k_2$ is a positive
integer degree. Along a path from the fixed vertex, the edge equation
forces the degrees to alternate $k_1,k_2,k_1,k_2,\ldots$.
Connectedness makes these the only degrees. If they are equal, the
graph is regular. If they differ, partition the vertices by these
two degrees. Each edge joins opposite parts by the edge equation;
both parts are nonempty because an edge exists. Thus this is a
bipartition and degrees are constant on each part. Counting each
edge at its endpoint in the first part gives $pk_1=|E(X)|$, and
counting in the second gives $qk_2=|E(X)|$.
\end{proof}
\lean{SimpleGraph.regular_or_semiregular_of_lineGraph_regular}
\status{New elementary proof. Isolated vertices must be removed;
they are invisible to the line graph.}
\used{[CGSS] Lemma 5.4 and Theorems 4.13, 5.6, 5.8.}

\begin{theorem}[Recovering a regular graph's spectrum]\label{thm:line-regular}
Let $X$ be a finite $k$-regular graph with at least one edge, and put
$n=|V(X)|$, $m=|E(X)|$. Then the spectrum of $L(X)$ determines $k,n$
and the full spectrum of $X$. In detail, with
$h=\lambda_{\max}(L(X))$ and $\rho=\rk(2I+A(L(X)))$,
\[
 k=(h+2)/2,\qquad n=2m/k,
\]
and $\Spec(X)$ is obtained by taking all nonzero eigenvalues of
$2I+A(L(X))$, adjoining $n-\rho$ zeros, and subtracting $k$ from every
entry. Equivalently, without negative exponents,
\[
 (t+2)^n\chi_{L(X)}(t)=(t+2)^m\chi_X(t+2-k).
\]
Thus two regular graphs with an edge and cospectral line
graphs are cospectral.
\end{theorem}
\begin{proof}
An edge and regularity force $k>0$. Every edge has degree $2k-2$ in
the line graph by Lemma~\ref{lem:line-connected}, whether or not
$X$ is connected. Theorem~\ref{thm:sp1} gives $h=2k-2$.
The multiset size of $L(X)$ is $m$, and the handshake equation
$nk=2m$ gives $n$. The incidence products are $kI+A(X)$ and
$2I+A(L(X))$ by Theorem~\ref{thm:incidence}. They have the same
nonzero eigenvalues with multiplicities by Theorem~\ref{thm:mn}.
The latter has rank $\rho$, so the former, of size $n$, has exactly
$n-\rho$ zeros. Shift by $-k$ to obtain the asserted spectrum.
Apply the characteristic-polynomial identity in
Theorem~\ref{thm:mn} to the $n\times m$ incidence matrix and then
substitute $t+2$ for its indeterminate. This gives the displayed
identity. All quantities in the reconstruction depend only on the
line-graph spectrum, proving the last assertion.
\end{proof}
\lean{SimpleGraph.spectrumMultiset_of_regular_lineGraph, SimpleGraph.cospectral_of_lineGraph_cospectral}
\status{New from M4, M6 and SP1, with full zero counts.}
\used{[CGSS] Lemma 5.5, Example 5.13 and C15. The edge hypothesis
excludes the counterexample of edgeless graphs of different orders.}

\begin{theorem}[Bipartite block spectra]\label{thm:block}
Let $B$ be a real $p\times q$ matrix of rank $\rho$, and let
$\lambda_1,\ldots,\lambda_\rho$ be the positive eigenvalues of $BB^T$,
with repetition. For real $s,t$, the matrix
\[
 C=\pmat{sI_p&B\\B^T&tI_q}
\]
has the two eigenvalues
\[
 \gamma_j^{\pm}=\frac{s+t\pm\sqrt{(s-t)^2+4\lambda_j}}2
 \quad(1\le j\le \rho),
\]
as well as $s$ repeated $p-\rho$ times and $t$ repeated $q-\rho$ times.
This is a multiset assertion, so coinciding displayed values add their
multiplicities. For $s=t=0$, it gives $\pm\sqrt{\lambda_j}$ and
$p+q-2\rho$ zeros. Furthermore
\[
 (C-sI)(C-tI)=\pmat{BB^T&0\\0&B^TB}.
\]
\end{theorem}
\begin{proof}
Choose orthonormal eigenvectors $\Zu_j$ of $BB^T$ for the positive
eigenvalues, and put $\Zv_j=B^T\Zu_j/\sqrt{\lambda_j}$.
Then
$\ip{\Zv_i}{\Zv_j}=\ip{\Zu_i}{BB^T\Zu_j}/
\sqrt{\lambda_i\lambda_j}$ is $1$ for $i=j$ and zero otherwise.
Moreover $B\Zv_j=\sqrt{\lambda_j}\Zu_j$ and
$B^T\Zu_j=\sqrt{\lambda_j}\Zv_j$.
The remaining $p-\rho$ orthonormal vectors in $\re^p$ lie in
$\ker B^T=\ker BB^T$. The remaining $q-\rho$ directions in $\re^q$
lie in $\ker B$: they are the orthogonal complement of
$\im B^T$, and the $\rho$ vectors $\Zv_j$ span that image by rank.
Thus the whole space is an orthogonal sum of the $\rho$ planes
spanned by $(\Zu_j,0),(0,\Zv_j)$ and these two kernel spaces.
On the $j$th plane the matrix of $C$ is
$\pmat{s&\sqrt{\lambda_j}\\\sqrt{\lambda_j}&t}$, with characteristic
polynomial $(z-s)(z-t)-\lambda_j$. Its discriminant is positive
because $\lambda_j>0$, so its roots are the two stated distinct
values. On the two kernel spaces $C$ acts by $s$ and $t$.
The dimensions add to $2\rho+(p-\rho)+(q-\rho)=p+q$, so this accounts for
every eigenvalue with multiplicity. Direct block multiplication of
$C^2-(s+t)C+stI$ gives the last displayed identity: the off-diagonal
blocks are $(s+t)B-(s+t)B=0$, and the diagonal blocks are the two
Gram products.
\end{proof}
\lean{Matrix.spectrumMultiset_fromBlocks_scalar, Matrix.fromBlocks_scalar_quadratic}
\status{New from M1 and the Gram-product facts.}
\used{[CGSS] Lemma 5.6, Remark 5.7 and Theorem 5.8; DS1--DS4.}

\begin{theorem}[Cospectral line graphs of bipartite graphs]\label{thm:line-bipartite}
Let $X_1$ be a connected finite bipartite graph with an edge, and
assume $L(X_1)$ is regular. Let $X_2$ be finite with no isolated
vertices and with $L(X_2)$ cospectral with $L(X_1)$. Then $X_2$ is
connected, and precisely one of the following alternatives holds:
\begin{enumerate}
\item $X_2$ is bipartite. Order the part sizes $p\le q$ in each
graph. The part sizes and respective degrees $k_1,k_2$ agree, and
the two biadjacency matrices $B_1,B_2$ have cospectral products
$B_1B_1^T,B_2B_2^T$.
\item $X_1$ is not regular; $X_2$ is regular and not bipartite;
$|V(X_1)|=|V(X_2)|+1$.
\end{enumerate}
\end{theorem}
\begin{proof}
The graph $L(X_1)$ is connected by Lemma~\ref{lem:line-connected}.
Its cospectral mate is regular and connected by
Corollary~\ref{cor:regular-mates}, and has the same positive number $m$ of vertices,
thus both root graphs have $m$ edges. Removing no vertices is needed
in $X_2$ because it has no isolates. Lemma~\ref{lem:line-connected}
gives its connectedness, and Theorem~\ref{thm:sp4} gives the
regular-or-semiregular alternative.
Put $h=\lambda_{\max}(L(X_1))$, $s=h+2$, and let $R$ be the common
rank of $2I+A(L(X_i))$, which is determined by the spectrum. In a
connected bipartite root graph, Theorem~\ref{thm:laplacian} and
Theorem~\ref{thm:incidence} show that its order is $N=R+1$ and that
its signless Laplacian has precisely one zero eigenvalue. Its full
spectrum is obtained from the common nonzero Gram spectrum by
adjoining that one zero.

For such a root graph, let the part sizes be $p\le q$ with degrees
$k_1\ge k_2>0$. Counting gives
\[
 p+q=N,\qquad k_1+k_2=s,\qquad pk_1=qk_2=m,
 \qquad k_1k_2=ms/N.
\]
Indeed $m/k_1+m/k_2=N$, so multiplication by $k_1k_2$ proves the
last equality. Thus the degrees are the ordered roots of
$z^2-sz+ms/N$, and then $p=m/k_1$, $q=m/k_2$.
They are determined, including the case of equal degrees.
The signless Laplacian is
$C=\pmat{k_1I&B\\B^T&k_2I}$. By Theorem~\ref{thm:block}, applying
$f(z)=(z-k_1)(z-k_2)$ gives $BB^T\oplus B^TB$.
Its spectrum is determined by $\Spec(C)$ and
Theorem~\ref{thm:calculus}. Every positive eigenvalue in this last
multiset occurs twice as often as in $BB^T$, by
Theorem~\ref{thm:mn}; halving these counts determines all positive
multiplicities of $BB^T$. Its remaining multiplicity at zero is $p$
minus their sum. This proves alternative (1) whenever $X_2$ is bipartite.

If $X_2$ is not bipartite, it must be regular, of degree $s/2$,
by Theorem~\ref{thm:sp4}. Its signless Laplacian has zero-dimensional
kernel, so its order is $R=N-1$. If $X_1$ were regular, its degree
would also be $s/2>0$. The two handshake equations with the common
edge count $m$ would then give equal orders, contradicting $N$ and
$N-1$. This proves all assertions of alternative (2). The alternatives
are exclusive because they assert opposite bipartiteness properties.
\end{proof}
\lean{SimpleGraph.lineGraph_cospectral_bipartite}
\status{New complete assembly of SP1--SP5 and block spectral mapping.}
\used{[CGSS] Lemma 5.6 and Theorem 5.8. The no-isolates hypothesis
on $X_2$ is necessary; otherwise any number of invisible isolates
could be adjoined.}

\begin{corollary}[The regular nonbipartite alternative]\label{cor:rival}
Let $X_1$ be connected bipartite semiregular with part sizes $p\le q$
and distinct degrees $k_1>k_2>0$, and let a connected regular
nonbipartite graph $X_2$ have a line graph cospectral with $L(X_1)$.
Put $\rho=\rk B$, $s=k_1+k_2$, where $B$ is the biadjacency matrix.
The positive eigenvalues of $BB^T$ include $k_1k_2$ once. Delete this
one occurrence and call the remaining list $\lambda_2,\ldots,\lambda_\rho$.
Then the spectrum of $X_2$ is
\[
 \ms{\frac s2,\quad
 \pm\sqrt{\lambda_j+\frac{(k_1-k_2)^2}{4}}\ (2\le j\le \rho),
 \quad\left(\frac{k_1-k_2}{2}\right)^{p-\rho},
 \quad\left(-\frac{k_1-k_2}{2}\right)^{q-\rho}}.
\]
\end{corollary}
\begin{proof}
The constant vectors give $B\Zj_q=k_1\Zj_p$ and
$B^T\Zj_p=k_2\Zj_q$, so $BB^T\Zj_p=k_1k_2\Zj_p$.
For $C=D(X_1)+A(X_1)$, the vector
$(k_1\Zj_p,k_2\Zj_q)$ is positive and satisfies $C\Zx=s\Zx$.
Its support is connected, so Theorem~\ref{thm:perron} and
Corollary~\ref{cor:positive-certificate} say that $s$ is simple.
Theorem~\ref{thm:block} pairs each occurrence of $k_1k_2$ with the
two roots $0,s$ of $(z-k_1)(z-k_2)=k_1k_2$. Thus $k_1k_2$ occurs
once. Also $C$ has exactly one zero by Theorem~\ref{thm:laplacian}.
The signless Laplacian of $X_2$ is nonsingular and has the nonzero
spectrum of $C$. Obtain its spectrum by deleting that one zero,
then subtract $s/2$ because $X_2$ has degree $s/2$ by
Theorem~\ref{thm:line-bipartite}. The paired roots for each remaining
$\lambda_j$ give the displayed square roots, the surviving $s$ gives
$s/2$, and the two kernel spaces of $B$ give the last two terms.
\end{proof}
\lean{SimpleGraph.lineGraph_cospectral_bipartite_spectrum}
\status{New, with explicit corrections to the endpoint and zero modes.}
\used{Error 14 of \code{cgss/errata.md}~\cite{errata} records both
endpoint and multiplicity corrections. [CGSS] Remark 5.7 and Theorem 5.8. The pair $0,s$ loses its
zero member before shifting; it does not give both $\pm s/2$ in
$X_2$. The last multiplicities sum to $p+q-2\rho$, which equals $q-p$
only when $B$ has full row rank. This states the needed hypothesis
missing from the final sentence of Remark 5.7.}

\section{Strongly regular graphs and two-eigenvalue matrices}\label{sec:srg}

\begin{theorem}[SR1: parameter and matrix identities]\label{thm:sr1}
Let $X$ be a finite strongly regular graph with parameters $(n,k,a,c)$:
it has $n$ vertices, is $k$-regular, and the numbers of common neighbours
of distinct adjacent and nonadjacent vertices are $a,c$, respectively.
Then
\[
 A^2=(k-c)I+(a-c)A+cJ,\qquad AJ=JA=kJ.
\]
If $n>0$, then $k(k-a-1)=c(n-k-1)$ as an equality of integers.
If $X$ is connected and $0<k<n-1$, then $0<c\le k$ and $a\le k-1$.
\end{theorem}
\begin{proof}
The diagonal of $A^2$ is $k$. Its off-diagonal entries count common
neighbours, and thus are $a$ or $c$ according as the pair is adjacent
or nonadjacent. These are exactly the entries of the first displayed
matrix. The row and column sums give $AJ=JA=kJ$.
Apply the first identity to $\Zj\ne0$ to obtain
$k^2=k-c+(a-c)k+cn=k+ak+c(n-k-1)$, which rearranges to the parameter
identity without truncated subtraction.
For an edge $uv$, common neighbours belong to $N(u)\setminus\{v\}$,
so $a\le k-1$. A nonadjacent pair exists by $k<n-1$, and its common
neighbours are a subset of a $k$-element neighbourhood, so $c\le k$.
For $c>0$, choose a shortest path between a nonadjacent pair in the
connected graph. The first and third vertices are not adjacent,
since that would shorten the path. They have its second vertex as
a common neighbour, forcing $c\ge1$.
\end{proof}
\lean{SimpleGraph.IsSRGWith.matrix_eq_integer, SimpleGraph.IsSRGWith.param_eq_integer}
\status{Mathlib: M7. New: cast-friendly ring forms with the stated
nontriviality bounds.}
\used{SR1; [CGSS] Theorem 4.13, Remark 4.14 and the spectral calculations of the examples.}

\begin{theorem}[SR3: spectrum and integral multiplicities]\label{thm:sr3}
Let $X$ be connected and strongly regular with parameters $(n,k,a,c)$,
where $0<k<n-1$. Put
\[
 \Delta=(a-c)^2+4(k-c),\qquad
 \theta_1=\frac{a-c+\sqrt\Delta}{2},\quad
 \theta_2=\frac{a-c-\sqrt\Delta}{2}.
\]
Then $\Delta>0$, $k>\theta_1>\theta_2$, and
\[
 \Spec(X)=\ms{k^1,\theta_1^{m_1},\theta_2^{m_2}},\qquad
 m_1=\frac{-k-(n-1)\theta_2}{\theta_1-\theta_2},\quad
 m_2=\frac{k+(n-1)\theta_1}{\theta_1-\theta_2}.
\]
Both displayed multiplicities are positive integers. Its spectral
idempotents are
\[
 E_0=\frac1nJ,\quad
 E_1=\frac{A-\theta_2I-\frac{k-\theta_2}{n}J}{\theta_1-\theta_2},\quad
 E_2=\frac{A-\theta_1I-\frac{k-\theta_1}{n}J}{\theta_2-\theta_1}.
\]
In particular $\rk E_i=\tr E_i=m_i$ for $i=1,2$,
$1+m_1+m_2=n$, and $k+m_1\theta_1+m_2\theta_2=0$.
\end{theorem}
\begin{proof}
By Theorem~\ref{thm:sr1}, $k-c\ge0$, so $\Delta\ge0$.
If it were zero, both its nonnegative terms would be zero, giving
$a=c=k$, contrary to $a\le k-1$. Thus $\Delta>0$ and the two real
roots are distinct. On $\Zj^\perp$, the matrix identity becomes
$A^2-(a-c)A-(k-c)I=0$. Apply it to each eigenvector in that invariant
space from Theorem~\ref{thm:ones}. Its eigenvalue is one of the
two displayed roots. The all-ones direction has eigenvalue $k$,
which is simple and largest by connectedness and
Theorem~\ref{thm:sp1}.
Both roots must occur on $\Zj^\perp$. Indeed, if just one value
$b$ occurred, then $A=bI+(k-b)J/n$, by its action on the
orthogonal decomposition $\spanop(\Zj)\oplus\Zj^\perp$.
Every off-diagonal entry would be equal. There is both an edge
($k>0$) and a distinct nonedge ($k<n-1$), with entries one and zero,
a contradiction. Since both occur and the $k$-eigenspace has
dimension one, both are strictly less than $k$.

Let their multiplicities be $m_1,m_2$, positive natural numbers by
occurrence and Theorem~\ref{thm:multiplicity}. Counting the basis
vectors gives $m_1+m_2=n-1$. Since $\tr A=0$,
$m_1\theta_1+m_2\theta_2=-k$ by Theorem~\ref{thm:walks}.
Substitute $m_2=n-1-m_1$ in the second equation and divide by the
nonzero number $\theta_1-\theta_2$ to obtain $m_1$; substituting
$m_1=n-1-m_2$ gives $m_2$. This proves integrality as a consequence
of equality with eigenspace dimensions, not as an assumed
feasibility condition.
The displayed $E_0$ is the identity on $\spanop(\Zj)$ and zero on
its orthogonal complement. For $E_1$, substitution on the $k$
direction gives zero; on a $\theta_1$ direction it gives one; on
a $\theta_2$ direction it gives zero. For $E_2$, the three
substitutions give zero, zero, one, respectively. Thus these are
exactly the spectral idempotents of Theorem~\ref{thm:calculus},
which supplies their trace, rank, products and spectral decomposition.
\end{proof}
\lean{SimpleGraph.IsSRGWith.spectrumMultiset_eq, SimpleGraph.IsSRGWith.spectralIdempotents, SimpleGraph.IsSRGWith.multiplicity_eq}
\status{New from M1, M7 and the idempotent algebra of
Theorem~\ref{thm:calculus}.}
\used{SR3, SR5's integrality filters, [CGSS] Examples 4.5--4.9 and Remark 4.14.}

\begin{theorem}[SR2: complements]\label{thm:sr2}
Let $X$ be connected strongly regular with $(n,k,a,c)$ and
$0<k<n-1$. Its complement has parameters
\[
 (n,n-k-1,n-2k+c-2,n-2k+a)
\]
and spectrum
\[
 \ms{(n-k-1)^1,(-1-\theta_1)^{m_1},(-1-\theta_2)^{m_2}},
\]
where the roots and multiplicities are those of
Theorem~\ref{thm:sr3}. Coincident values in this list must be combined;
the complement need not be connected.
\end{theorem}
\begin{proof}
Complement degree is $n-k-1$. If $u,v$ are nonadjacent in $X$,
the union of their original neighbourhoods has size $2k-c$ and
contains neither $u$ nor $v$. Among the other $n-2$ vertices,
exactly $n-2-2k+c$ lie in neither neighbourhood; these are the
common neighbours in the complement. If $u,v$ are adjacent in
$X$, their neighbourhood union has size $2k-a$ and contains
both $u,v$. Removing them leaves $2k-a-2$ forbidden vertices
among $n-2$, so the common-neighbour count is $n-2k+a$.
These counts are nonnegative integers and prove the parameter
formula. The spectral formula is Theorem~\ref{thm:ones}, because
$A(\comp X)=J-I-A(X)$. It gives multiplicities as a multiset even
when the new all-ones eigenvalue also occurs on $\Zj^\perp$.
\end{proof}
\lean{SimpleGraph.IsSRGWith.spectrumMultiset_compl}
\status{Mathlib: the complement parameters in M7. New: the spectral
formula including possible collisions.}
\used{SR2, GQ3, [CGSS] Remark 4.14 and the complement conventions for the examples.}

\begin{corollary}[SR4: least eigenvalue minus two]\label{cor:sr4}
Let $X$ be connected strongly regular with $(n,k,a,c)$,
$0<k<n-1$, and least eigenvalue $-2$. Then
\[
 \theta_1=a-c+2\in\ints,\quad k=c+2\theta_1,\quad
 m_1=\frac{2n-k-2}{\theta_1+2},\quad
 \rk(2I+A)=1+m_1.
\]
In particular the positive integer $\theta_1+2$ divides
$2n-k-2$ in $\ints$, and the rank is one plus that quotient.
\end{corollary}
\begin{proof}
The least root of Theorem~\ref{thm:sr3} is $\theta_2=-2$.
The sum and product of the two roots of
$z^2-(a-c)z-(k-c)$ give
$\theta_1-2=a-c$ and $-2\theta_1=c-k$.
These yield the first two equations and integrality of $\theta_1$.
The multiplicity formula in that theorem yields
$m_1=(2(n-1)-k)/(\theta_1+2)$.
We have $\theta_1>-2$, and $k>0$, so in $2I+A$ the all-ones
eigenvalue $k+2$ and the $m_1$ copies of $\theta_1+2$ are positive,
while the remaining eigenvalues vanish. This proves the rank formula.
The integer equation $(\theta_1+2)m_1=2n-k-2$ gives divisibility.
\end{proof}
\lean{SimpleGraph.IsSRGWith.parameters_of_leastEigenvalue_eq_neg_two, SimpleGraph.IsSRGWith.rank_two_add_adjMatrix}
\status{New from SR3 and the rank bridge.}
\used{SR4--SR5, the parameter certificate C18, and [CGSS] Remark 4.14.}

\begin{definition}[Designs and their graphs]\label{def:design}
A $2$-design $D$ consists of finite point and block types and an incidence
relation, with parameters $(v,k,b,r,\lambda)$: $v>1$ points, $b>0$
labelled blocks, $k>1$ points in each block, $r$ blocks at each point,
and $\lambda$ blocks at each pair of distinct points. Repeated blocks
are allowed. In Lean the fields are \code{v k b r lam}.
Its zero-one point-block matrix is $N=D.\mathrm{incMatrix}$.
The graph $I(D)=D.\mathrm{incidenceGraph}$ has the disjoint union of the
point and block types as vertices and exactly the incidences as edges.
The flag graph is $F(D)=D.\mathrm{flagGraph}=L(I(D))$.
\lean{Design, Design.incMatrix, Design.incidenceGraph, Design.flagGraph}
\status{Proposed definitions following \code{cgss/notation.md}, Section 7.}
\end{definition}

\begin{lemma}[A distinguished constant eigendirection]\label{lem:two-eigenvalues}
Let $G$ be real symmetric of order $n>0$, with $G\Zj=b\Zj$, and
suppose every eigenvalue of its restriction to $\Zj^\perp$ equals $s$
(the restriction may have dimension zero). Then
\[
 G=sI+\frac{b-s}{n}J.
\]
\end{lemma}
\begin{proof}
The proposed matrix and $G$ act by $b$ on $\Zj$ and by $s$ on the
orthonormal basis of $\Zj^\perp$ from Theorem~\ref{thm:ones}.
Equality on this basis proves the identity.
\end{proof}
\lean{Matrix.IsHermitian.eq_smul_one_add_smul_allOnes_of_two_eigenvalues}
\status{New scalar-action consequence of the all-ones decomposition.}
\used{The design Gram characterization and SR3.}

\begin{theorem}[Designs from two Gram eigenvalues]\label{thm:two-eigenvalues}
Let $N$ be a $v\times b$ zero-one matrix, where $v>1$,
with constant positive row sum $r$ and column sum $k>1$. Let
$\lambda>0$ satisfy $r(k-1)=\lambda(v-1)$. If
$\Spec(NN^T)=\ms{(rk)^1,(r-\lambda)^{v-1}}$, then
\[
 NN^T=(r-\lambda)I+\lambda J.
\]
Thus any two distinct rows have inner product $\lambda$, and $N$
is the incidence matrix of a $2$-design with these parameters, with
columns treated as block labels. If $r>\lambda$, then $b\ge v$.
\end{theorem}
\begin{proof}
For $NN^T$, the row and column sums give
$NN^T\Zj_v=rk\Zj_v$. Here $rk>r-\lambda$ because
$rk-(r-\lambda)=r(k-1)+\lambda=\lambda v>0$.
The spectrum hypothesis thus makes the $rk$-eigenspace
one-dimensional, and it is spanned by $\Zj_v$; all remaining
eigendirections have eigenvalue $r-\lambda$. Lemma~\ref{lem:two-eigenvalues}
gives coefficient $(rk-r+\lambda)/v=\lambda$ on $J$.
The entries of $NN^T$ count common incident columns, proving the
design assertion without requiring distinct columns.
If $r>\lambda$, both $r-\lambda$ and $rk$ are positive, so
$\rk(NN^T)=v$. By Theorem~\ref{thm:gram}, $\rk N=v\le b$.
\end{proof}
\lean{Matrix.gram_eq_of_two_eigenvalues}
\status{New from M1 and the all-ones decomposition.}
\used{DS1, DS3--DS4 and the step ``it is the incidence matrix of a
2-design'' in [CGSS] Theorem 5.8.}

\begin{theorem}[Spectra of design incidence and flag graphs]\label{thm:design-spectra}
Let $N=D.\mathrm{incMatrix}$ be the zero-one incidence matrix of a $2$-design $D$ on $v>1$
points and $b$ labelled blocks of size $k>1$, with replication $r$
and pair count $\lambda>0$. Suppose $r>\lambda$. Put $m=vr=bk$.
Let $I(D)$ be its bipartite incidence graph and $F(D)=L(I(D))$ its
flag graph. Define
\[
 \gamma^{\pm}=\frac{r+k\pm\sqrt{(r-k)^2+4(r-\lambda)}}2.
\]
Then $b\ge v$, $m-v-b+1\ge0$, and
\[
 \begin{aligned}
 \Spec(I(D))&=\ms{\sqrt{rk},-\sqrt{rk},
    (\sqrt{r-\lambda})^{v-1},(-\sqrt{r-\lambda})^{v-1},0^{b-v}},\\
 \Spec(F(D))&=\ms{r+k-2,(\gamma^+-2)^{v-1},
   (\gamma^--2)^{v-1},(k-2)^{b-v},(-2)^{m-v-b+1}}.
 \end{aligned}
\]
Repeated displayed values are combined. Both graphs are connected,
and the flag graph is $(r+k-2)$-regular.
\end{theorem}
\begin{proof}
Counting incidences by points and blocks gives $vr=bk$.
Fix a point and count pairs consisting of a block through it and
a different point of that block. The two counts give
$r(k-1)=\lambda(v-1)$. Entries count common blocks, so
$NN^T=(r-\lambda)I+\lambda J$.
Corollary~\ref{cor:scalar-ones} gives its eigenvalues
$r-\lambda+\lambda v=rk$ once and $r-\lambda$ with multiplicity
$v-1$. Both are positive, and their difference is $\lambda v>0$.
Thus $N$ has rank $v$ by Theorem~\ref{thm:gram}, and $b\ge v$.
Apply Theorem~\ref{thm:block} with the two scalar blocks zero to
the adjacency matrix of $I(D)$ to obtain its displayed spectrum.

Any two distinct points lie in a common block because
$\lambda>0$, so a path of length two joins them. Every block has
a point because $k>0$. These paths join all vertices of $I(D)$;
there is an edge because $v,r>0$. Thus both graphs are connected
by Lemma~\ref{lem:line-connected}. That lemma also gives the flag
degree $r+k-2$. The signless Laplacian of $I(D)$ has scalar blocks
$rI_v,kI_b$. The positive Gram eigenvalue $rk$ gives the two
roots $0,r+k$ in Theorem~\ref{thm:block}; each occurrence of
$r-\lambda$ gives $\gamma^+,\gamma^-$, and the $b-v$ kernel
directions give $k$. Both $\gamma$ values are positive: their
sum is $r+k>0$, their product is
$rk-r+\lambda=\lambda v>0$, and they are real. Thus zero has
exactly one occurrence. Theorem~\ref{thm:incidence} and
Theorem~\ref{thm:mn} say that $2I+A(F(D))$ has this nonzero
spectrum and $m-v-b+1$ zeros. This also proves that the last
exponent is nonnegative. Subtracting two proves the flag spectrum.
\end{proof}
\lean{Design.spectrumMultiset_incidenceGraph, Design.spectrumMultiset_flagGraph}
\status{New specialization of the proved block and incidence
identities. Repeated blocks cause no difficulty because block
labels are distinct vertices.}
\used{DS1, DS3--DS4, [CGSS] Theorem 5.8 and Corollaries 5.10--5.11.
For $(v,k,b,r,\lambda)=(6,3,10,5,2)$ this gives C15's list.}

\begin{corollary}[The trivial design and complete bipartite graphs]\label{cor:complete-bipartite}
For $p,q\ge1$,
\[
 \Spec(L(K_{p,q}))=
 \ms{p+q-2,(p-2)^{q-1},(q-2)^{p-1},(-2)^{(p-1)(q-1)}}.
\]
Values that coincide are combined. This covers the design case
$r=\lambda>0$, in which every point belongs to every block.
\end{corollary}
\begin{proof}
The biadjacency matrix $B$ is all ones. Its image is the span of
$\Zj_p$, of dimension one, and $BB^T=qJ_p$ has positive
eigenvalue $pq$ once. The signless Laplacian has scalar blocks
$qI_p,pI_q$. Theorem~\ref{thm:block} gives eigenvalues
$0,p+q$, together with $q$ repeated $p-1$ times and $p$ repeated
$q-1$ times. The graph is connected: opposite parts are joined
by an edge and same-part vertices by a path through the other
nonempty part. Its edge incidence matrix thus has rank
$p+q-1$ by Theorem~\ref{thm:incidence}. The $pq$-dimensional
other Gram product has the preceding nonzero spectrum and
$pq-p-q+1=(p-1)(q-1)$ zeros, by Theorem~\ref{thm:mn}.
Subtracting two gives the displayed formula. In a design with
$r=\lambda>0$ on $v$ points, the count $r(k-1)=\lambda(v-1)$ gives $k=v$.
Each block thus contains every point. This makes its incidence
graph complete bipartite, as claimed.
\end{proof}
\lean{SimpleGraph.spectrumMultiset_lineGraph_completeBipartite}
\status{New specialization retaining collisions and the cases
$p=1$ or $q=1$.}
\used{[CGSS] Corollary 5.9 and the trivial-design qualification
after Theorem 5.8.}

\begin{theorem}[The polynomial on the constant direction]\label{thm:hoffman-poly}
Let $X$ be a connected finite $k$-regular graph on $n>0$ vertices,
and let $f$ be a real polynomial which vanishes at every eigenvalue
other than $k$. Then
\[
 f(A)=\frac{f(k)}n J.
\]
Conversely this matrix identity implies that $f$ vanishes at every
eigenvalue on $\Zj^\perp$.
\end{theorem}
\begin{proof}
The $k$-eigenspace is exactly $\spanop(\Zj)$ by
Theorem~\ref{thm:sp1}. In the eigenbasis of
Theorem~\ref{thm:ones}, $f(A)$ acts by $f(k)$ on that direction
and by zero on the others, precisely as the stated multiple of
$J$. Conversely apply the identity to an eigenvector in
$\Zj^\perp$; the right side is zero and the left side is
$f(\mu)\Zv$. Since $\Zv\ne0$, $f(\mu)=0$.
\end{proof}
\lean{SimpleGraph.IsRegularOfDegree.aeval_eq_smul_allOnes_of_connected}
\status{New consequence of spectral idempotents.}
\used{[CGSS] Theorem 5.8's final matrix identity and HR2:
for spectrum $4^1,2^3,0^3,(-2)^5$, take $f(z)=z^3-4z$ to get
$A^3=4A+4J$.}

\section{Smith graphs}\label{sec:smith}

We will give SM2 and SM3 explicitly, so that the use of
Perron--Frobenius in Theorem 3.8 has no combinatorial black box.
\begin{definition}[Smith models]\label{def:smith-models}
Write $T(p,q,r)$ for the tree consisting of one vertex of degree three
and three internally disjoint arms of positive edge-lengths $p,q,r$.
Define $\widetilde{\mathrm A}_{\ell}$ to be the cycle on $\ell+1$
vertices, $\ell\ge2$; $\widetilde{\mathrm D}_4=K_{1,4}$; and, for
$\ell\ge5$, define $\widetilde{\mathrm D}_{\ell}$ by a path of length
$\ell-4$ with two new leaves at each endpoint. Define
\[
 \widetilde{\mathrm E}_6=T(2,2,2),\quad
 \widetilde{\mathrm E}_7=T(1,3,3),\quad
 \widetilde{\mathrm E}_8=T(1,2,5).
\]
These are the Smith graphs of [CGSS], Figure 1; the finite-type and
affine-type names are kept distinct.
A labelled realization uses a central vertex and arm vertices $(a,j)$,
$1\le j\le h_a$, joined consecutively and to the centre at $j=1$.
For $\widetilde{\mathrm D}_\ell$ label the path consecutively and label
each of its four added leaves separately; cycles use cyclic indices.
Write $\code{SimpleGraph.smithE}(h)$ for $h\in\{6,7,8\}$.
\lean{SimpleGraph.threeArmTree, SimpleGraph.smithA, SimpleGraph.smithD, SimpleGraph.smithE}
\status{Proposed finite graph definitions with the displayed parameter domains.}
\end{definition}

\begin{theorem}[SM2: positive integer eigenvectors]\label{thm:smith-vectors}
Every Smith graph has largest eigenvalue $2$, of multiplicity one.
For a cycle, a certificate vector is $\Zj$. For
$\widetilde{\mathrm D}_{\ell}$, it has value $2$ on the central path
(or on the centre when $\ell=4$) and value $1$ on the four leaves.
For $T(p,q,r)$ in the three exceptional cases, take central value
$L=3,4,6$, respectively; at distance $j$ on an arm of length $h$ take
\[
 \Zx(j)=\frac{L(h+1-j)}{h+1}\qquad(1\le j\le h).
\]
For a Smith graph on $N$ vertices, $2I-A$ is positive semidefinite
of rank $N-1$. Deleting any vertex from a Gram representation with
this Gram matrix leaves a basis of its span.
\end{theorem}
\begin{proof}
All the stated coordinates are positive integers. On a cycle each
vertex has two neighbours of value one. On an affine $D$ graph,
a leaf sees value two; an internal path vertex sees two values two;
an endpoint of a nontrivial central path sees values $2,1,1$;
and in $\widetilde{\mathrm D}_4$ the centre sees four values one.
In each case the neighbour sum is twice the coordinate.
On an exceptional arm the coordinates decrease in an arithmetic
progression from $L$ to $L/(h+1)$, with next value zero beyond the
leaf. Thus the sum of the two adjacent values is twice the current
one; at the leaf the only actual neighbour has twice its value.
At the centre the neighbour sum is
$L(3-1/(p+1)-1/(q+1)-1/(r+1))=2L$, because the reciprocals sum to
one for $(2,2,2),(1,3,3),(1,2,5)$. This proves $A\Zx=2\Zx$.
Corollary~\ref{cor:positive-certificate} makes $2$ the largest
eigenvalue; connectedness and Theorem~\ref{thm:perron} make it simple.
Thus $2I-A$ is positive semidefinite and has precisely one zero
eigenvalue. A graph with one vertex deleted is a nonempty proper
subgraph, so has largest eigenvalue strictly less than $2$ by
Theorem~\ref{thm:strict}. Its $2I-A$ is positive definite by
Theorem~\ref{thm:psd}. This is the principal Gram submatrix of the
remaining $N-1$ vectors, so they are independent by
Theorem~\ref{thm:gram}. The original span has dimension $N-1$,
so they form a basis.
\end{proof}
\lean{SimpleGraph.largestEigenvalue_smithA, SimpleGraph.largestEigenvalue_smithD, SimpleGraph.largestEigenvalue_smithE}
\status{New parametric certificates and their proof from SM1.}
\used{SM2, C17, [CGSS] Theorem 3.8 and the sentence following Figure 1.}

\begin{lemma}[Trees of maximum degree two]\label{lem:tree-path}
A finite tree whose vertices have degree at most two is a nonempty path.
\end{lemma}
\begin{proof}
Choose a longest simple path, which exists by finiteness and nonemptiness.
An outside vertex has a path meeting it for the first time. Attachment
at an endpoint extends the chosen path; attachment internally gives
three distinct neighbours. Both are impossible. Thus every vertex is
on the chosen path; any additional edge would create a cycle.
\end{proof}
\lean{SimpleGraph.IsTree.isPath_of_degree_le_two}
\status{New finite-tree lemma, reusable on the components after deletion.}
\used{The two shape cases in Smith's theorem.}

\begin{theorem}[SM3: Smith's theorem]\label{thm:smith}
A connected finite graph has largest eigenvalue $2$ iff it is
isomorphic to a Smith graph. It has largest eigenvalue less than
$2$ iff it is a path, a tree $T(1,1,r)$ for $r\ge1$, or one of
$T(1,2,2),T(1,2,3),T(1,2,4)$.
\end{theorem}
\begin{proof}
First suppose a connected graph contains a Smith graph as a subgraph.
If it is exactly that graph, its largest eigenvalue is $2$ by
Theorem~\ref{thm:smith-vectors}. If the subgraph is proper, strict
monotonicity makes its largest eigenvalue greater than $2$.
We will classify connected graphs containing no Smith subgraph.
Such a graph has no cycle, so is a tree; it has no vertex of degree
at least four, since four of the incident edges would give $K_{1,4}$.
It cannot have two vertices of degree three. Indeed, choose two
whose distance is minimal. Internal vertices on their unique path
have degree two; otherwise one would be a closer degree-three
vertex. At each endpoint choose its two neighbours off this path.
These four neighbours are distinct and off the path, since a
coincidence would give a cycle in the tree. The path and these four
incident edges form an affine $D$ subgraph.

If the tree has no degree-three vertex, it is a path, including a
single vertex. To see this directly, take a longest simple path.
A vertex outside it would join it by a unique path; attachment at
an endpoint extends it and attachment internally gives degree at
least three, both impossible. If there is one degree-three vertex,
delete it. Each of the three components has degree at most two
away from the attachment, thus is a path by the same longest-path
argument. The tree is thus $T(p,q,r)$, with
$1\le p\le q\le r$ after reordering. If $p\ge2$, truncate all
arms to length two to obtain $\widetilde{\mathrm E}_6$.
Thus $p=1$. If $q\ge3$, truncate to $(1,3,3)$ to obtain
$\widetilde{\mathrm E}_7$. Thus $q=1$ or $q=2$.
When $q=2$, avoidance of $(1,2,5)$ gives $2\le r\le4$.
This proves the asserted list of graphs avoiding Smith subgraphs.

Every graph on that list is a proper subgraph of a Smith graph.
A path on $N$ vertices embeds into a cycle on $N+2$ vertices.
For $T(1,1,r)$, adjoin two new leaves at the far endpoint of the
arm of length $r$. This gives the affine $D$ tree with central
path length $r$. The three remaining trees embed properly in
$T(1,2,5)$ by extending the last arm. Strict monotonicity now
makes every listed graph's largest eigenvalue less than $2$.
It also shows that a graph with largest eigenvalue at most $2$
cannot contain a proper Smith subgraph. Combining this with the
exhaustive avoidance argument proves both directions of both
classification statements.
\end{proof}
\lean{SimpleGraph.smith, SimpleGraph.largestEigenvalue_lt_two_iff}
\status{New complete combinatorial proof from SM1 and SM2.}
\used{SM3 and [CGSS] Theorem 3.8.}

\section{Exact certificates and rational transfer}\label{sec:certificates}

\begin{lemma}[Moments on an annihilating support]\label{lem:annihilating-moments}
Let $A$ be real symmetric of finite order and let $(\theta_i)_{i<s}$
be distinct real nodes, where $s\ge0$. If
$\prod_{i<s}(A-\theta_iI)=0$, every eigenvalue belongs to the nodes and
\[
 \tr A^j=\sum_{i<s}\mult_A(\theta_i)\theta_i^j\qquad(j\ge0).
\]
\end{lemma}
\begin{proof}
Apply the annihilator to a nonzero eigenvector. The scalar product
$\prod_{i<s}(\mu-\theta_i)$ vanishes, so a factor vanishes.
Theorem~\ref{thm:trace} sums powers over all eigenvalues; partition
that sum by its unique node, using distinctness, to obtain the formula.
If $s=0$, the annihilator is $I=0$. This forces the index set to be
empty, and every asserted sum and trace is zero.
\end{proof}
\lean{Matrix.IsHermitian.eigenvalues_mem_of_annihilating, Matrix.IsHermitian.trace_pow_eq_multiplicities}
\status{Implemented in \code{LineGraphs/Spectral/Basic.lean}.}
\used{Both matrices in a fixed-support certificate.}

\begin{theorem}[Moments determine multiplicities on a known support]\label{thm:certificate}
Let $A$ be real symmetric of order $n$, let
$\theta_0,\ldots,\theta_{s-1}$ be distinct real numbers with $s\ge0$,
and let $m_0,\ldots,m_{s-1}$ be nonnegative integers. Suppose
\[
 \prod_{i=0}^{s-1}(A-\theta_iI)=0,\qquad
 \tr A^j=\sum_{i=0}^{s-1}m_i\theta_i^j\quad(0\le j<s).
\]
Then
$\Spec(A)=\ms{\theta_0^{m_0},\ldots,\theta_{s-1}^{m_{s-1}}}$.
In particular $\sum_i m_i=n$, including the empty node list.
\end{theorem}
\begin{proof}
Put $d_i=\mult_A(\theta_i)$. Lemma~\ref{lem:annihilating-moments}
and the assumed moments give $\sum_i(d_i-m_i)\theta_i^j=0$ for $j<s$.
Apply \code{Matrix.eq_zero_of_forall_pow_sum_mul_pow_eq_zero}
to the injective node function and the vector of real differences.
It gives $d_i=m_i$; injectivity of natural-number casts gives equality
in $\nn$. The support assertion of that lemma and these counts give
the multiset by extensionality. Taking its cardinality and applying
Theorem~\ref{thm:multiplicity} gives $\sum_i m_i=n$, also for $s=0$.
For $s>0$ this is also the zeroth moment, with $0^0=1$.
\end{proof}
\lean{Matrix.IsHermitian.spectrumMultiset_eq_of_annihilatingPolynomial_trace}
\status{Implemented in \code{LineGraphs/Spectral/Basic.lean}, using M4
Vandermonde uniqueness, including its empty system.}
\used{The kernel certificates of C13 and C15, and any fixed-support spectrum computation.}

\begin{corollary}[Two matrices on a common annihilating support]\label{cor:certificate-pair}
Let $A,B$ be real symmetric matrices on possibly different finite index
sets. Suppose they annihilate the product on the same distinct nodes
$(\theta_i)_{i<s}$, $s\ge0$, and have equal traces in degrees $0\le j<s$.
Then $\Spec(A)=\Spec(B)$ and $\chi_A=\chi_B$.
\end{corollary}
\begin{proof}
Lemma~\ref{lem:annihilating-moments} for $B$ supplies moments with
$m_i=\mult_B(\theta_i)$. The trace equalities give these same moments
for $A$, so Theorem~\ref{thm:certificate} applied to $A$ gives precisely
$B$'s multiset. The support statement of the lemma for $B$ identifies
that multiset; Proposition~\ref{prop:cospectral} gives polynomial equality.
This includes $s=0$, when both annihilators force empty index sets.
\end{proof}
\lean{Matrix.IsHermitian.charpoly_eq_of_annihilatingPolynomial_trace}
\status{Implemented in \code{LineGraphs/Spectral/Basic.lean}.}
\used{C13 and C15 certificates across different matrix index types.}

\begin{theorem}[Similarity and intertwining certificates]\label{thm:similarity}
Let $A,B,Q$ be square matrices of the same order over a field, with
$\det Q\ne0$ and $QA=BQ$. Then $\chi_A=\chi_B$.
In particular $Q^2=I$ and $QAQ=B$ suffice.
For integer matrices $A,B,Q,R$ and an integer $d\ne0$, it is enough
to check
\[
 QA=BQ,\qquad QR=RQ=dI.
\]
These identities imply equality of the real characteristic polynomials
of $A$ and $B$. For adjacency matrices, they prove graph cospectrality.
\end{theorem}
\begin{proof}
An invertible $Q$ has inverse $Q^{-1}$, and $B=QAQ^{-1}$.
Thus
$tI-B=Q(tI-A)Q^{-1}$; determinant multiplicativity gives
$\chi_B=\det Q\,\chi_A\,\det Q^{-1}=\chi_A$.
When $Q^2=I$, its inverse is $Q$ and the conjugation hypothesis
is exactly the required equation. For the integer certificate,
cast the identities into $\rats$; $(1/d)R$ is an explicit two-sided
inverse of $Q$. The field assertion gives equality of rational
characteristic polynomials. The inclusion $\rats\hookrightarrow\re$
preserves determinant: its finite permutation sum uses only products,
sums and signs, all preserved by a ring map. Thus the equality
persists over $\re$. The graph assertion is
Definition~\ref{def:graph}.
\end{proof}
\lean{Matrix.charpoly_eq_of_invertible_intertwiner, Matrix.charpoly_eq_of_scaled_inverse_intertwiner, SimpleGraph.cospectral_of_intertwiner}
\status{Mathlib: \code{Matrix.charpoly_units_conj} and
\code{Matrix.charpoly_map} (M4). Both matrix intertwiner wrappers are implemented in
\code{LineGraphs/Spectral/Similarity.lean}, using
\code{Matrix.charpoly_mul_comm}; the graph wrapper is proposed here.}
\used{The kernel certificate of C15, and switching. An explicit inverse
certificate avoids computing a large symbolic determinant.}

\begin{theorem}[Seidel switching]\label{thm:seidel}
Let $X$ be a finite simple graph and $U$ a subset of its vertices.
Define $X^U$ by complementing adjacency across $U$ and its complement
and retaining adjacency inside the two parts. Let $T$ be diagonal,
with entry $-1$ on $U$ and $1$ elsewhere. Then
\[
 S(X^U)=T S(X)T,\qquad T^2=I.
\]
Thus Seidel switching preserves the Seidel spectrum with multiplicity.
The matrix used in [CGSS] Corollary 5.10 is $-S=2A-J+I$, and its
spectrum is obtained from that of $S$ by negating every entry.
\end{theorem}
\begin{proof}
Both sides have diagonal zero. For distinct vertices in the same
part the sign product is one, and neither the adjacency nor the
Seidel entry changes. Across the parts the product is $-1$ and
adjacency is complemented, changing the Seidel entry between $1$
and $-1$. This proves the conjugation identity entry by entry.
The diagonal signs square to one. Apply
Theorem~\ref{thm:similarity} and Proposition~\ref{prop:cospectral}
to the real symmetric Seidel matrices. Finally use
Theorem~\ref{thm:calculus} with the polynomial $-z$.
\end{proof}
\lean{SimpleGraph.seidelMatrix_switch, SimpleGraph.spectrumMultiset_seidelMatrix_switch}
\status{New Seidel definitions and entry calculation; Mathlib and
Theorem~\ref{thm:similarity} give the conjugation step.}
\used{SW2, [CGSS] Section 4 switching examples and Corollary 5.10, C13.}

\begin{theorem}[Rational eigenvalues of integer matrices]\label{thm:rational}
If an integer square matrix has a rational eigenvalue $b$, viewed as
an eigenvalue of its real matrix, then $b$ is an integer.
\end{theorem}
\begin{proof}
The characteristic polynomial $f$ is monic with integer coefficients,
and $f(b)=0$: the real eigenvector makes the real determinant zero,
and injectivity of the rational inclusion makes the rational
evaluation zero. Its degree is positive, since a constant monic
polynomial has no root. Write $b=h/d$ in lowest terms, with
$d>0$ and $\gcd(h,d)=1$. If the degree is $N\ge1$, multiply
$f(h/d)=0$ by $d^N$. The leading term is $h^N$ and every other
term is divisible by $d$. Thus $d\mid h^N$. Coprimality of $h$
and $d$ implies coprimality of $h^N$ and $d$: repeatedly apply
Euclid's lemma, or Bezout's identity, to the factors $h$.
The only positive divisor of $h^N$ coprime to $h^N$ is $1$,
so $d=1$ and $b=h\in\ints$.
\end{proof}
\lean{Matrix.eigenvalue_isInteger_of_rat}
\status{Mathlib: the monic rational-root theorem and
\code{Matrix.charpoly_map} (M4). New: the matrix-eigenvalue wrapper.}
\used{LA6, SR4's arithmetic context, and $x\mid n$ in [CGSS] Section 6.}

\begin{lemma}[Adjacency matrix guards]\label{lem:adjacency-guards}
Let $A$ be a symmetric integer matrix on a finite set, with zero diagonal
and entries in $\{0,1\}$. The relation $A(u,v)=1$ defines a simple
graph whose integer adjacency matrix is $A$.
\end{lemma}
\begin{proof}
The three guards are the three fields of Mathlib's
\code{Matrix.IsAdjMatrix}. Symmetry of $A$ gives symmetry of the
relation, and the zero diagonal gives irreflexivity; this graph is
\code{Matrix.IsAdjMatrix.toGraph}. On the diagonal both adjacency and
input entries are zero. Off the diagonal, the zero-one guard makes the
test for equality with one return exactly the input entry; this is
\code{Matrix.IsAdjMatrix.adjMatrix_toGraph_eq}.
\end{proof}
\lean{Matrix.IsAdjMatrix.of_symmetric_zero_one}
\status{Mathlib: \code{Matrix.IsAdjMatrix} and its graph (M6); the proposed
name only packages the three guards as its constructor. The integer-list
version is \code{Computations.graphOfMatrix} with
\code{Computations.adjMatrix_graphOfMatrix} in
\code{LineGraphs/Computations/GraphCertificates.lean}.}
\used{Each proposed adjacency matrix, separately, in a graph certificate.}

\begin{theorem}[A precise Boolean certificate interface]\label{thm:checker}
Encode an integer $n\times n$ matrix by a list of $n$ rows of length
$n$. Consider the following three checkers, using exact integers and
natural numbers, with no floating-point operations.
\begin{enumerate}
\item A spectral checker first verifies the square shape and symmetry,
then a possibly empty list of distinct integer nodes $\theta_i$ and an
equally long list of natural multiplicities $m_i$. It verifies the
matrix product $\prod_i(A-\theta_iI)=0$ and the trace identities
of Theorem~\ref{thm:certificate} for $0\le j<s$.
\item A largest-eigenvalue checker verifies square shape, $n>0$,
symmetry, nonnegative entries, a length-$n$ integer vector with
strictly positive entries, and $A\Zx=b\Zx$ for the proposed integer $b$.
\item A similarity checker verifies that $A,B,Q,R$ have the same
square shape, that $d\ne0$, and the three identities in
Theorem~\ref{thm:similarity}.
\end{enumerate}
Returning \code{true} implies, respectively, the stated real multiset
spectrum, the stated largest eigenvalue, or equality of real
characteristic polynomials. Graph certificates additionally require zero
diagonal, zero-one entries, and symmetry for every proposed adjacency
matrix. In the similarity checker these guards apply to both $A$ and $B$.
\end{theorem}
\begin{proof}
We will specify the arithmetic so that the reflection claim has a
precise meaning. For equal-length lists define their dot product
recursively by
$\operatorname{dot}([],[])=0$ and
$\operatorname{dot}(a::v,b::w)=ab+\operatorname{dot}(v,w)$.
Matrix multiplication takes the dot product of each row with each
column, after shape checks. Addition and scalar multiplication act
entrywise, the identity has entries $1$ on equal indices and $0$
otherwise, and trace is the sum of diagonal entries. A list equality
test compares lengths and then corresponding entries. A universal
test is the conjunction of its finitely many Boolean tests, with
empty conjunction true.
Induction on the common length proves that dot computes
$\sum_i v(i)w(i)$: the empty sum is zero and the successor step
splits off the first summand. Applying this to each row and column
proves correctness of matrix multiplication. Induction on list
length proves that a universal test returns true exactly when all
its tests do, and the equality test does so exactly when all entries
agree. Thus true shape tests give well-defined matrices on
$\operatorname{Fin}(n)$, and true algebraic tests give the stated
identities over $\ints$.

For the polynomial product start at $I$ and multiply successively
by $A-\theta_iI$; induction on the processed nodes shows that the
accumulator is that initial segment's product. For powers start at
$P_0=I$ and put $P_{j+1}=P_jA$; induction gives $P_j=A^j$.
The trace test thus checks exactly the specified moments,
including $j=0$. The inclusion $\ints\hookrightarrow\re$ preserves
all finite sums and products, so these matrix and trace identities
remain true over $\re$. It also preserves inequalities and is
injective; the node distinctness, vector positivity and nonzero
denominator tests thus give the same hypotheses over $\re$.
Apply Theorem~\ref{thm:certificate},
Corollary~\ref{cor:positive-certificate}, or
Theorem~\ref{thm:similarity}, respectively.
With the graph guards, Lemma~\ref{lem:adjacency-guards} defines
adjacency on $\operatorname{Fin}(n)$ by $A(u,v)=1$ and shows that its
adjacency matrix is exactly the encoded matrix. In the similarity
checker it applies to $A$ and to $B$ separately. This proves each graph
conclusion.
\end{proof}
\lean{Matrix.spectrumCertificate_sound, Matrix.positiveEigenvectorCertificate_sound, Matrix.similarityCertificate_sound}
\status{Implemented as \code{Computations.checkSpectrum_sound},
\code{Computations.checkPositiveEigenvector_sound}, and
\code{Computations.checkSimilarity_sound} in
\code{LineGraphs/Computations/SpectrumChecker.lean},
\code{PositiveEigenvector.lean}, and \code{SimilarityChecker.lean},
respectively, with the three displayed \code{Matrix} names as aliases.
The graph similarity corollary
\code{Computations.cospectral_graphOfMatrix_of_checkSimilarity}
requires \code{checkAdjacency} for both matrices. A Python run is not a kernel proof.}
\used{The kernel certificates of C13, C15 and C17. All finite checking
can be proved with kernel reduction after these soundness lemmas.}

The concrete spectral inputs needed for C13 and C15 are small.
The column of moments below starts at degree zero; the annihilating
polynomial is the product of $z-\theta$ over the listed nodes.
\begin{center}
\small
\begin{tabular}{@{}llll@{}}
\toprule
Computation & Nodes & Multiplicities & Moments\\
\midrule
C13, $A$ & $4,2,0,-2$ & $1,3,3,5$ & $12,0,48,48$\\
C13, $S$ & $-5,-1,3$ & $3,3,6$ & $12,0,132$\\
C13, $-S$ & $5,1,-3$ & $3,3,6$ & $12,0,132$\\
C15 & $6,4,1,0,-2$ & $1,5,4,5,15$ & $30,0,180,420,2820$\\
\bottomrule
\end{tabular}
\end{center}
The sign convention is $S=J-I-2A$, as in SW2 of \code{cgss/prerequisites.md}~\cite{plan};
Corollary 5.10 of [CGSS] uses $2A-J+I=-S$. The C13 rows apply to both sets of twelve vectors printed in [CGSS].
The C15 row applies to $L(L(\mathrm{Petersen}))$ and to the flag graph
of the $2$-$(6,3,2)$ design. The Python checks construct both
matrices in each case and check these products and moments exactly.
They also check the positive integer vectors of
Theorem~\ref{thm:smith-vectors} for C17. The entries of this table are
certificate data; the preceding soundness theorem is the proof that
checking them suffices. Each instantiated certificate requires a Lean proof
that its specific checker returns true.

\section{Further spectral applications used in CGSS}\label{sec:applications}

\begin{theorem}[Scaled idempotents and singular principal minors]\label{thm:scaled}
Let $C$ be real symmetric of order $n$ and let $b>0$ satisfy
$C^2=bC$. Then $C\succeq0$, its eigenvalues lie in $\{0,b\}$,
and $\rk C=\tr C/b$. If $B$ is an $n\times q$ matrix with $CB=0$,
then $\rk C\le n-\rk B$. Every principal submatrix of order
greater than $\rk C$ is singular.
\end{theorem}
\begin{proof}
Applying the equation to an eigenvector gives
$\mu(\mu-b)=0$, so each eigenvalue is $0$ or $b>0$.
Theorems~\ref{thm:psd} and \ref{thm:trace} give positivity and
$\tr C=b\rk C$. The equation $CB=0$ gives
$\im B\sbs\ker C$. Taking dimensions and using rank-nullity
gives $\rk B\le n-\rk C$. A principal submatrix is a composite
of the inclusion of a coordinate subspace, the map $C$, and a
coordinate projection, so its rank is at most $\rk C$. If its
order is larger, it has nonzero kernel by rank-nullity and
thus zero determinant.
\end{proof}
\lean{Matrix.IsHermitian.posSemidef_of_sq_eq_pos_smul, Matrix.IsHermitian.rank_eq_trace_div_of_sq_eq_smul, Matrix.det_principal_submatrix_eq_zero_of_rank_lt}
\status{New short interfaces from M1--M4.}
\used{HC3: $C^2=35C$, $CN^T=0$, $n=21$, $\rk N=15$ give positivity
and rank at most $6$ for Connor's matrix.}

\begin{theorem}[Symmetric matrices and rank-one Gram matrices]\label{thm:symmetric-space}
The real symmetric $\ell\times\ell$ matrices form an inner product
space for $\ip{P}{Q}=\tr(PQ)$, of dimension $\ell(\ell+1)/2$.
For $\ell>0$ the traceless subspace has dimension
$\ell(\ell+1)/2-1$. For arbitrary real vectors $\Zv_u$, their rank-one
matrices $P_u=\Zv_u\Zv_u^T$ have Gram matrix
$\tr(P_uP_v)=\ip{\Zv_u}{\Zv_v}^2$.
Thus, if $n$ unit vectors in $\re^\ell$ have distinct-pair
inner products in $\{0,1/2\}$, then $n\le\ell(\ell+1)/2$.
\end{theorem}
\begin{proof}
For symmetric $P,Q$,
$\tr(PQ)=\sum_{u,v}P(u,v)Q(u,v)$. This is bilinear and symmetric,
and $\tr(P^2)=\sum_{u,v}P(u,v)^2$ is nonnegative, with equality
only when every entry is zero. Thus it is an inner product.
Take one matrix with a single diagonal entry one for each index,
and one with entries one at $(u,v),(v,u)$ for each unordered
distinct pair. These matrices have disjoint entry supports; their
coefficients in a linear combination are read from the indicated
entries, so they are independent. Every symmetric matrix is the
sum with coefficients its corresponding entries, so they span.
Their number is $\ell+\ell(\ell-1)/2=\ell(\ell+1)/2$.
For $\ell>0$, the trace linear map onto $\re$ is surjective since
a single diagonal basis matrix has trace one. Rank-nullity gives
the traceless dimension.

Expansion of entries gives
\[
 \tr(P_uP_v)
 =\sum_{i,j}\Zv_u(i)\Zv_u(j)\Zv_v(j)\Zv_v(i)
 =\Bigl(\sum_i\Zv_u(i)\Zv_v(i)\Bigr)^2.
\]
In the last assertion let $G=I+A/2$ be the original Gram matrix,
where $A$ records the pairs with inner product $1/2$.
The rank-one Gram matrix is $I+A/4=(I+G)/2$.
Its form on any nonzero $\Zx$ is
$(\norm{\Zx}^2+\Zx^TG\Zx)/2>0$, because $G\succeq0$.
It thus has rank $n$. By Theorem~\ref{thm:gram}, the $P_u$
span an $n$-dimensional subspace of the symmetric matrices, proving
the dimension bound. For $n=0$ the bound is immediate, and if
$\ell=0$ no unit vector exists, so this also handles that case.
\end{proof}
\lean{Matrix.finrank_symmetric, Matrix.finrank_traceless_symmetric, Matrix.gram_rankOne_eq_inner_sq, SimpleGraph.card_le_symm_dimension_of_representation}
\status{New elementary dimension and rank-one calculation;
Mathlib supplies \code{Sym2.card} (M8), finite-dimensional inner
product and rank facts.}
\used{LA5 and [CGSS] Theorem 4.3's bound $n\le36$ in dimension $8$;
the centred version underlies Theorem 4.4.}

\begin{definition}[Hadamard pair matrices]\label{def:hadamard-pairs}
For a real matrix $H$ on a finite row type $U$ and column type $V$,
let $P=\code{Matrix.hadamardPairs}\ H$ have rows indexed by the
unordered two-element subsets $\{i,j\}\subseteq U$, with
$P(\{i,j\},u)=H(i,u)H(j,u)$. This does not depend on the ordering
of $i,j$. For square $H$ of order $n$ and real $x>0$, define
$\code{Matrix.hadamardSignedAdj}\ H\ x=(PP^T-nI)/x$ and write it $C_H$.
The name does not by itself assert that the entries are $0,1,-1$;
that requires an additional hypothesis.
\lean{Matrix.hadamardPairs, Matrix.hadamardSignedAdj}
\status{Proposed definitions following \code{cgss/notation.md}, Section 7.}
\end{definition}

\begin{theorem}[The Hadamard pair matrix]\label{thm:hadamard}
Let $H$ be an $n\times n$ real matrix with entries $\pm1$,
$HH^T=nI$, and $n\ge2$. Let $P$ have one row for each unordered
pair of distinct rows of $H$, with
$P(\{i,j\},u)=H(i,u)H(j,u)$, and put $m=\binom n2$. Then
\[
 P^TP=\frac{n^2}{2}I-\frac n2J,\qquad
 \Spec(PP^T)=\ms{(n^2/2)^{n-1},0^{m-n+1}},\qquad
 \rk P=n-1.
\]
For a positive real $x$, put $C_H=(PP^T-nI)/x$. Then
\[
 (xC_H+nI)\left(xC_H-\frac{n(n-2)}2I\right)=0,
\]
and its spectrum consists of $n(n-2)/(2x)$ with multiplicity
$n-1$ and $-n/x$ with multiplicity $m-n+1$.
If $n>2$, $x$ is a positive integer, and $C_H$ has integer
entries, then $x\mid n$.
\end{theorem}
\begin{proof}
Since $n>0$, the identity $H(H^T/n)=I$ makes $H$ surjective and
thus invertible in finite equal dimension. Its inverse is $H^T/n$,
so also $H^TH=nI$. For two column indices $u,v$, put
$b_i=H(i,u)H(i,v)$. The identity
$2\sum_{i<j}b_ib_j=(\sum_i b_i)^2-\sum_i b_i^2$ follows by
separating diagonal terms in the square of the sum. Here
$\sum_i b_i=n$ if $u=v$, and zero otherwise, by column
orthogonality; each $b_i^2=1$. Thus
$(P^TP)(u,v)=n^2\delta_{uv}/2-n/2$, as asserted.
On $\Zj$, this matrix acts by zero. On $\Zj^\perp$ it acts by
$n^2/2>0$. Its rank is $n-1$. Theorem~\ref{thm:mn} transfers its
nonzero spectrum to $PP^T$, which has $m-n+1$ zeros and rank $n-1$.
This number is $(n-1)(n-2)/2\ge0$.
Polynomial calculus proves $PP^T(PP^T-(n^2/2)I)=0$;
substituting $PP^T=xC_H+nI$ gives the displayed identity and
the shifted, scaled eigenvalue list. For $n>2$ the zero multiplicity
of $PP^T$ is positive, so $-n/x$ really occurs as a rational
eigenvalue of the integer matrix $C_H$. By
Theorem~\ref{thm:rational}, $-n/x=z\in\ints$. Multiplying by
$x$ gives $n=x(-z)$, which is $x\mid n$.
\end{proof}
\lean{Matrix.IsHadamard.hadamardPairs_transpose_mul, Matrix.IsHadamard.spectrumMultiset_hadamardPairs_gram, Matrix.IsHadamard.dvd_order_of_integral_hadamardSignedAdj}
\status{New assembly of entry counting, M3--M4, and polynomial
calculus over Mathlib's \code{Matrix.IsHadamard} (M8).}
\used{[CGSS] Section 6, C16 and HD1. This gives the spectral and
divisibility input to Theorem 6.2; the root-system classification
and four-row agreement counting are separate results.}

\begin{remark}[A row-sum assertion in Section 6]\label{rem:hadamard-row}
The additional printed identity $xC_HJ=-nJ$ does not follow from
the stated definition of $P$. It is equivalent to
$PP^T\Zj_m=0$, and thus to $P^T\Zj_m=0$ by
Theorem~\ref{thm:gram}. The valid identity $P\Zj_n=0$ instead
sums the products of \emph{two rows} of $H$, which are orthogonal;
these are different assertions, on different spaces.
For an explicit counterexample, take the normalized Sylvester matrix
\[
 H=\pmat{1&1&1&1\\1&-1&1&-1\\1&1&-1&-1\\1&-1&-1&1},\qquad x=4.
\]
Its six pair rows are the three nonconstant rows of $H$, each twice.
Thus $PP^T$ is, after reordering, three blocks $4J_2$.
The matrix $C_H=(PP^T-4I)/4$ is three blocks
$\pmat{0&1\\1&0}$, whose row sums are $1$, not $-1$.
This example also satisfies Hypothesis 6.1 with $n=x=4$:
the unique four-row set agrees in exactly one column, which is
$(n+x)/8$. Error 9 in \code{cgss/errata.md}~\cite{errata} records this
row-sum error. Its repair also admits the shorter argument: for $n>2$,
the eigenvalue $-n/x$ has multiplicity $(n-1)(n-2)/2>0$, so the
monic rational-root theorem gives $x\mid n$. The divisibility conclusion in
Theorem~\ref{thm:hadamard} is thus deliberately proved from
the actual eigenvalue $-n/x$, without this row-sum assertion.
\end{remark}

\section{Further directions}\label{sec:further}

The matrix multiplicity bridge, polynomial trace and multiset
spectral mapping, Rayleigh equality, and rectangular Gram-spectrum
lemmas are natural additions to Mathlib. They use existing
definitions and would shorten several independent developments.
The signless Laplacian kernel and its bipartite-component basis
fit beside the existing Laplacian file. Symmetric Perron--Frobenius
and interlacing can follow after the Rayleigh interface; a natural
question is which support-graph formulation gives the most useful
bridge to Mathlib's existing \code{Matrix.IsIrreducible}.

For graphs, we will keep the multiset definition separate from the
set spectrum, and keep graph cospectrality separate from vertex
cospectrality. The strongly regular eigenvalue and idempotent
formulas are also suitable for upstreaming once their nontriviality
hypotheses are reflected in the API. We will retain the standing
positive-semidefinite hypothesis for the root-system development:
only statements requiring an extreme eigenvalue need a nonempty
vertex type.

The companion structured document supplies named proof steps.
The spectral certificate, positive-vector and similarity lemmas
already have proofs in \code{LineGraphs/Spectral/}; that directory also
defines matrix and graph switching equivalence. The Boolean soundness
lemmas are implemented in \code{LineGraphs/Computations/}, while the
Seidel switching theorems and
other proposed interfaces remain mathematical specifications. The checks in \code{checks/check_spectraltools.py}
use integer products for C13, C15 and C17, and NumPy and SymPy for
independent small sanity checks. They do not establish Lean kernel
proofs or the non-spectral classification claims of CGSS.
No ordinary proof stated as a theorem here is left open. The
unqualified statements that required repair were the empty or
isolated-vertex cases, the Seidel sign in the first statement of C13, the endpoint
and rank qualifications in Remark 5.7, and the row-sum assertion
discussed in Remark~\ref{rem:hadamard-row}.

\begin{thebibliography}{9}
\bibitem{errata} \emph{Errata to CGSS}, \code{cgss/errata.md}, Errors 9 and 14.
\bibitem{CGSS}
P. J. Cameron, J. M. Goethals, J. J. Seidel and E. E. Shult,
\emph{Line graphs, root systems, and elliptic geometry},
J. Algebra \textbf{43} (1976), 305--327.
We read it with the corrections recorded in \code{cgss/errata.md}.
\bibitem{mathlib}
The Mathlib community, \emph{Mathlib}, commit \code{082e2d37}
(4 September 2026), the version pinned in \code{lake-manifest.json}.
Exact declarations and files are listed in Section~\ref{sec:catalogue}.
\bibitem{plan}
\emph{Prerequisites}, \code{cgss/prerequisites.md}; and
\emph{Notation}, \code{cgss/notation.md}, Section 10,
in this repository.
\end{thebibliography}
\end{document}
